% =====================================================================
%  numerical-scheme.tex
%
%  The majorant-frequency Monte Carlo engine behind BF_v_no_resampling_v2.py:
%  the baseline algorithm, the ten tricks it is built from, and the
%  measurement layer that turns a run into a number.
%
%  TWO WAYS TO USE THIS FILE
%  -------------------------
%  (a) Standalone: compile as it is.  It is self-contained and cites nothing
%      external, so there are no undefined references.
%  (b) As a Part of main.tex: delete everything between the two lines marked
%      "PREAMBLE -- CUT FROM HERE" and "PREAMBLE -- CUT TO HERE", delete the
%      final \end{document}, save under Parts/, and add
%          % =====================================================================
%  numerical-scheme.tex
%
%  The majorant-frequency Monte Carlo engine behind BF_v_no_resampling_v2.py:
%  the baseline algorithm, the ten tricks it is built from, and the
%  measurement layer that turns a run into a number.
%
%  TWO WAYS TO USE THIS FILE
%  -------------------------
%  (a) Standalone: compile as it is.  It is self-contained and cites nothing
%      external, so there are no undefined references.
%  (b) As a Part of main.tex: delete everything between the two lines marked
%      "PREAMBLE -- CUT FROM HERE" and "PREAMBLE -- CUT TO HERE", delete the
%      final \end{document}, save under Parts/, and add
%          % =====================================================================
%  numerical-scheme.tex
%
%  The majorant-frequency Monte Carlo engine behind BF_v_no_resampling_v2.py:
%  the baseline algorithm, the ten tricks it is built from, and the
%  measurement layer that turns a run into a number.
%
%  TWO WAYS TO USE THIS FILE
%  -------------------------
%  (a) Standalone: compile as it is.  It is self-contained and cites nothing
%      external, so there are no undefined references.
%  (b) As a Part of main.tex: delete everything between the two lines marked
%      "PREAMBLE -- CUT FROM HERE" and "PREAMBLE -- CUT TO HERE", delete the
%      final \end{document}, save under Parts/, and add
%          % =====================================================================
%  numerical-scheme.tex
%
%  The majorant-frequency Monte Carlo engine behind BF_v_no_resampling_v2.py:
%  the baseline algorithm, the ten tricks it is built from, and the
%  measurement layer that turns a run into a number.
%
%  TWO WAYS TO USE THIS FILE
%  -------------------------
%  (a) Standalone: compile as it is.  It is self-contained and cites nothing
%      external, so there are no undefined references.
%  (b) As a Part of main.tex: delete everything between the two lines marked
%      "PREAMBLE -- CUT FROM HERE" and "PREAMBLE -- CUT TO HERE", delete the
%      final \end{document}, save under Parts/, and add
%          \include{Parts/numerical-scheme}
%      next to the existing \include{Parts/unified-cascade_v4}.
%      main.tex already loads amsmath, amssymb, graphicx and hyperref; you
%      must additionally add to its preamble:
%          \usepackage{algorithm}
%          \usepackage{algpseudocode}
%          \usepackage{booktabs}
% =====================================================================

% ---------------------------------------------- PREAMBLE -- CUT FROM HERE
\documentclass[a4paper,12pt]{article}

\usepackage[T1]{fontenc}
\usepackage[utf8]{inputenc}
\usepackage[english]{babel}
\usepackage{amsmath,amssymb}
\usepackage{geometry}
\geometry{margin=2.5cm}
\usepackage{graphicx}
\usepackage{booktabs}
\usepackage{algorithm}
\usepackage{algpseudocode}
\usepackage[colorlinks=true,linkcolor=blue,citecolor=blue,urlcolor=blue]{hyperref}

\newcommand{\minj}{m_{\mathrm{inj}}}
\newcommand{\msink}{m_{\mathrm{sink}}}
\newcommand{\fmaj}{f^{\mathrm{maj}}}
\newcommand{\code}[1]{\texttt{#1}}

\begin{document}

\title{The numerical scheme: majorant-frequency Monte Carlo\\ for the mass cascade}
\author{}
\date{}
\maketitle
% ---------------------------------------------- PREAMBLE -- CUT TO HERE

\section{Scope and design principle}
\label{sec:num_scope}

The engine integrates the stochastic mass cascade in four configurations,
\[
  \text{process} \in \{\text{coagulation},\ \text{fragmentation}\},
  \qquad
  \text{system} \in \{\text{closed},\ \text{open}\},
\]
through a \emph{single} code path. Only two things differ between the four
cases:

\begin{enumerate}
  \item what an accepted event does to the two selected particles, and
  \item the boundary conditions --- injection plus an absorbing sink, or
        neither.
\end{enumerate}

Everything else --- majorant construction, pair selection, the clock, the
bookkeeping --- is shared. This is a deliberate constraint, not a
convenience: a discrepancy between two of the four cases cannot then be
blamed on two different integrators.

The variant documented here carries no particle weights. One simulated
particle is one physical particle, $w \equiv 1$ for the whole run, so
$\sum_i w_i$ and $\sum_i w_i m_i$ are conserved to machine precision at
every step. The measured mass drift is identically zero, and any nonzero
value is a bug rather than a statistical excursion. Section~\ref{sec:num_cost}
states what that costs.

\subsection{State}

The state is four parallel arrays over live particles, plus a binning of
mass space:

\begin{center}
\small
\begin{tabular}{@{}ll@{}}
\toprule
\code{mass[i]}       & particle mass $m_i$ \\
\code{inj\_time[i]}  & birth time, used for the isochrones (Sec.~\ref{sec:trick7}) \\
\code{bin\_of[i]}    & index of the mass bin holding $i$ \\
\code{pos\_in\_bin[i]} & position of $i$ inside that bin's index list \\
\midrule
\code{edges[0..B]}   & $B+1$ logarithmically spaced mass-bin edges \\
\code{bins[b]}       & list of particle indices in bin $b$ \\
\code{N\_bin[b]}     & occupancy $N_b$ \\
\bottomrule
\end{tabular}
\end{center}

A particle of mass $m$ belongs to bin
$b(m) = \max\bigl(0,\min(B-1,\ \mathrm{searchsorted}(\text{edges},m)-1)\bigr)$.
The clamping in that expression is a trap and is discussed in
Sec.~\ref{sec:num_pitfalls}.

\section{The baseline algorithm}
\label{sec:num_algorithm}

\subsection{Setup}

\begin{algorithm}[h]
\caption{Setup: bins, majorant table, rate vector}
\label{alg:setup}
\begin{algorithmic}[1]
\State bin every initial particle; form the occupancies $N_b$
\State \textbf{majorant table:} for all $b,c$ \quad
       $\fmaj_{bc} \gets K(\text{edges}[b+1],\ \text{edges}[c+1])$
       \Comment{upper corners}
\State $T_b \gets \sum_c \fmaj_{bc} N_c$
\State $\text{row}_b \gets N_b\,(T_b - \fmaj_{bb})$
\State $R \gets \sum_b \text{row}_b$
\end{algorithmic}
\end{algorithm}

The corner rule $\fmaj_{bc}=K(\text{edges}[b+1],\text{edges}[c+1])$ bounds
$K$ over the whole bin pair \emph{provided $K$ is non-decreasing in both
arguments}. That is the single structural assumption on the kernel, and it
is verified at run time rather than assumed: every accepted trial checks
$K(m_i,m_j) \le \fmaj_{bc}$ and raises if it fails.

By construction
\begin{equation}
  R \;=\; \sum_b N_b\bigl(T_b - \fmaj_{bb}\bigr)
    \;=\; \sum_{i \neq j} \fmaj_{ij}
    \;=\; 2 \sum_{i<j} \fmaj_{ij},
  \label{eq:num_R}
\end{equation}
i.e.\ $R$ runs over \emph{ordered} pairs. The factor two in
Eq.~\eqref{eq:num_R} is the origin of the factor two in the clock,
Eq.~\eqref{eq:num_clock}.

\subsection{Main loop}

\begin{algorithm}[h]
\caption{Main loop: one majorant trial}
\label{alg:mainloop}
\begin{algorithmic}[1]
\While{no stopping condition is met}
  \State $\Delta t \gets 2V/(wR)$; \quad $t \gets t + \Delta t$
         \Comment{the clock, Sec.~\ref{sec:trick5}}
  \State \textbf{injection (open systems):}
         $\nu \gets \nu + q\,\Delta t / w$; \quad
         $k \gets \lfloor \nu \rfloor$; \quad $\nu \gets \nu - k$;
         add $k$ particles of mass $\minj$
  \State draw a row bin $b$ with $P(b) \propto \text{row}_b$
  \State draw a column bin $c$ with
         $P(c) \propto N_c \fmaj_{bc}$, using $(N_b-1)\fmaj_{bb}$ for $c=b$
  \State draw $i$ uniformly from \code{bins[b]}, $j \neq i$ uniformly from \code{bins[c]}
  \State \textbf{thin:} accept with probability $K(m_i,m_j)/\fmaj_{bc}$;
         otherwise \textbf{continue} \Comment{null event}
  \State apply the event (Alg.~\ref{alg:event}); apply the absorbing sink
  \State record a snapshot if the cadence says so (Sec.~\ref{sec:trick10})
\EndWhile
\end{algorithmic}
\end{algorithm}

A rejected trial is not an error: it is the thinning that makes the
dominating process exact. The acceptance ratio, reported as \code{acc} in
the run log, is $O(1)$ by construction of the bin-pair table --- measured
$0.95$ for the geometric kernel on bins of $0.1$~dex, where the corner
overshoot is $10^{0.1\lambda}$ per argument.

\subsection{The event}

\begin{algorithm}[h]
\caption{The accepted event --- the only place the four cases differ}
\label{alg:event}
\begin{algorithmic}[1]
\If{coagulation}
  \State $m_{\mathrm{new}} \gets m_i + m_j$
  \State inherit the age of $i$ from $(i,j)$ by \code{age\_rule} (Sec.~\ref{sec:trick7})
  \State $m_i \gets m_{\mathrm{new}}$; rebin $i$; \textbf{delete} $j$
  \State \textbf{if} open \textbf{and} $m_{\mathrm{new}} \ge \msink$ \textbf{then}
         delete $i$; $M_{\mathrm{out}} \mathrel{+}= w\,m_{\mathrm{new}}$;
         $n_{\mathrm{out}} \mathrel{+}= 1$
\Else \Comment{fragmentation}
  \State $p \gets \arg\max(m_i,m_j)$; \quad $s \gets \arg\min(m_i,m_j)$
  \State \textbf{if} $m_s < f\,m_p$ \textbf{then continue}
         \Comment{disruption threshold, Sec.~\ref{sec:trick8}}
  \State $\xi \sim U(0,1)$; \quad $m_a \gets \xi m_p$, \ $m_b \gets (1-\xi)m_p$
  \State both fragments inherit $t_{\mathrm{born}}$ of the single parent $p$
  \State $m_p \gets m_a$; rebin $p$; \textbf{append} a particle of mass $m_b$
  \State \textbf{if} open: delete every fragment with $m \le \msink$,
         highest index first
\EndIf
\end{algorithmic}
\end{algorithm}

Note the sign asymmetry, which governs everything in
Sec.~\ref{sec:num_pitfalls}: coagulation is $2 \to 1$ and therefore a
\emph{sink} of particle number, fragmentation is $1 \to 2$ and therefore a
\emph{source}.

\section{The ten tricks}
\label{sec:num_tricks}

\subsection{Majorant collision frequency}
\label{sec:trick1}

Computing the exact pair rate for every pair is $O(N^2)$ per step. Instead a
majorant $\fmaj \ge K$ drives a dominating Poisson process and events are
thinned with acceptance $K/\fmaj$. The thinned process is exact: the
majorant affects the cost, never the statistics.

\emph{Sources.} Ivanov \& Rogasinsky, Sov.\ J.\ Numer.\ Anal.\ Math.\
Modelling \textbf{3}, 453 (1988) [MCF scheme in DSMC]; Garcia, van den
Broeck, Aertsens \& Serneels, Physica~A \textbf{143}, 535 (1987) [majorant
for coagulation]; Gillespie, J.~Atmos.\ Sci.\ \textbf{32}, 1977 (1975)
[exact stochastic coalescence].

\subsection{Bin-pair majorant table}
\label{sec:trick2}

The majorant is a matrix $\fmaj_{bc}$ over pairs of mass bins, not a single
global number. What that buys is \emph{cost} and \emph{checkability}. It is
worth being precise about which, because a majorant cannot buy correctness:
it cancels out of the physical rate exactly.

\paragraph{The majorant cancels.}
Ticks of the dominating process arrive at rate $wR/(2V)$; a tick selects the
ordered pair $(i,j)$ with probability $\fmaj_{ij}/R$; it is accepted with
probability $K_{ij}/\fmaj_{ij}$. The rate of accepted events on that pair is
therefore
\begin{equation}
  \frac{wR}{2V}\cdot\frac{\fmaj_{ij}}{R}\cdot\frac{K_{ij}}{\fmaj_{ij}}
  \;=\; \frac{w\,K_{ij}}{2V},
  \label{eq:num_cancel}
\end{equation}
and summing over unordered pairs gives $w\sum_{i<j}K_{ij}/V$ --- the
Smoluchowski rate, independent of $\fmaj$ altogether. The majorant cancels
twice, once in the selection weight and once in the acceptance. \emph{Any}
valid majorant therefore returns the same physics; a bad one only wastes
trials. Equation~\eqref{eq:num_cancel} holds because the clock is advanced on
every \emph{trial} rather than on every event --- see
Sec.~\ref{sec:trick5}, which spells out the one edit that breaks it.

\paragraph{What the table actually buys.}
A single global $f_{\max} = K(m_{\max},m_{\max})$ is a legitimate majorant and
gives identical statistics. It is nevertheless the wrong choice, for three
reasons:

\begin{itemize}
  \item \textbf{Acceptance.} A typical pair is accepted with probability
        $\langle K\rangle/f_{\max} \simeq (m_0/m_{\max})^{\lambda}$. Across
        $D$ decades that is $10^{-\lambda D}$: for the geometric kernel over
        four decades, $2\times10^{-3}$, i.e.\ some five hundred trials per
        event against $1.05$ with the table. The table instead keeps the
        acceptance $O(1)$ in \emph{every} bin pair, because the bound is
        evaluated at the corners of the pair actually drawn; the overshoot is
        then only $10^{0.1\lambda}$ per argument at $0.1$~dex bins.
  \item \textbf{The stall guard misfires.} \code{max\_stall\_tries} counts
        trials without an accepted event. At an acceptance of
        $2\times10^{-3}$ its default of $2\times10^{6}$ is only
        $\sim\!4000$ events' worth of honest work, so a healthy run can be
        killed with the diagnosis ``acceptance collapsed''.
  \item \textbf{A lagging maximum is a genuine bias.} A \emph{running} global
        maximum that is not refreshed before some particle grows past it
        gives $\fmaj < K$, and the thinning is then silently wrong --- this,
        unlike the two points above, is a correctness failure. The corner
        table is static, and every trial checks $K \le \fmaj$ explicitly, so
        the same situation raises instead of biasing.
\end{itemize}

\emph{Source.} Eibeck \& Wagner, SIAM J.~Sci.\ Comput.\ \textbf{22}, 802
(2000),
\url{https://doi.org/10.1137/S1064827599353488}.

\subsection{$O(1)$ bin membership}
\label{sec:trick3}

\code{bins[b]} is a list of particle indices and \code{pos\_in\_bin[i]} is
the position of $i$ inside it. Removal swaps the last element into the hole,
so insertion and deletion are $O(1)$ with no reallocation. The same
swap-and-pop idiom is used on the particle arrays themselves when a particle
is deleted, which is why every index held across a deletion must be
revalidated.

\subsection{Incremental rate maintenance}
\label{sec:trick4}

When one bin count changes by $\delta$,
\[
  T \mathrel{+}= \delta\, \fmaj_{\cdot b},
  \qquad
  \text{row} \gets N \odot (T - \mathrm{diag}\,\fmaj),
  \qquad
  R \gets \textstyle\sum_b \text{row}_b ,
\]
an $O(B)$ vector operation instead of recomputing the full matrix--vector
product. Selection uses a vectorised \code{cumsum} plus \code{searchsorted}
rather than a Python loop; that loop is otherwise the dominant cost. For
very large $B$ a Fenwick tree gives $O(\log B)$ (Goodson \& Kraft,
J.~Comput.\ Phys.\ \textbf{183}, 210 (2002)).

\subsection{The clock}
\label{sec:trick5}

The step is
\begin{equation}
  \Delta t \;=\; \frac{2V}{w\,R},
  \label{eq:num_clock}
\end{equation}
and three separate things about it are easy to get wrong.

\paragraph{Ordered versus unordered pairs.}
$R$ counts every unordered pair twice, Eq.~\eqref{eq:num_R}, while the
physical majorant rate runs over unordered pairs. Hence the numerator is
$2V$, not $V$. Pair \emph{selection} is unaffected --- both orderings lead to
the same unordered pair and the factors cancel --- so the error is invisible
in the dynamics and rescales $t$ by exactly two, breaking every comparison
with an analytic solution.

\paragraph{Weights.}
Here $w$ is a constant, but it is still written into the step so that the
formula matches the weighted variant exactly and a constant $w \neq 1$ simply
rescales the density. Nothing in this variant ever changes $w$, so the clock
cannot be corrupted by weight bookkeeping. That is the entire reason the
baseline exists.

\paragraph{Quadratic scaling.}
$R \propto N^2$ because the process is binary, so
\begin{equation}
  t \;\simeq\; \frac{2E}{N^2}
  \label{eq:num_t_of_E}
\end{equation}
after $E$ trials. A step linear in $N$ describes a unary process and gives
the wrong $N(t)$. Equation~\eqref{eq:num_t_of_E} has a practical corollary
that is easy to meet by accident: \emph{raising $N$ at fixed event budget
makes the run shorter in physical time, not longer}. Any gate expressed as an
absolute time and tuned at one $N$ is never crossed at $10N$.

\paragraph{Time advances on every trial, not on every event.}
\code{t\_phys} is incremented immediately after $\Delta t$ is formed, some
fifty lines before the acceptance test, so null (thinned) trials move the
clock too. That is precisely what makes the cancellation
\eqref{eq:num_cancel} work, and moving the increment inside the accepted
branch is the one edit that turns the majorant from a cost into a
\emph{bias}. Physical time would then advance by $2V/(wR)$ per accepted
event, so the reported event rate would be $R$ instead of
$R\,\langle K/\fmaj\rangle = \sum K$, and with a global $f_{\max}$,
\[
  \Delta t_{\rm eff} \;\propto\; \frac{1}{N^2 f_{\max}}
                     \;\propto\; \frac{1}{N^2 m_{\max}^{\lambda}} .
\]
The kernel would then enter through the \emph{extreme} of the distribution
instead of through the typical mass. In an open system $m_{\max}$ is pinned
at the sink and does not grow, so the $m_0$ dependence leaves the clock
entirely, the drift degenerates to $dm_0/dt \propto m_0$ (exponential,
$\beta = 1$), and the constant-flux argument then returns $\alpha = -2$ for
every kernel. This is a third, purely numerical route to the same $-2$
fingerprint that Sec.~\ref{sec:trick8} produces physically.

\subsection{Deterministic residual injection}
\label{sec:trick6}

Injection uses a fractional accumulator,
\[
  \nu \mathrel{+}= q\,\Delta t/w,
  \qquad k = \lfloor \nu \rfloor,
  \qquad \nu \mathrel{-}= k,
\]
rather than a Poisson draw. This removes shot noise at the injection scale.
It is a modelling choice, and it is recorded in the output metadata rather
than left implicit.

\subsection{Isochrone bookkeeping}
\label{sec:trick7}

Every particle carries its birth time; its age is $\tau = t - t_{\rm born}$.
Snapshots accumulate a two-dimensional histogram in $(\tau, m)$, whose rows
are the isochrones and whose sum over rows is the time-averaged spectrum.

For \textbf{fragmentation} the age of a fragment is unambiguous: fragments
descend from exactly one parent. For \textbf{coagulation} a merged particle
has two ancestors and an inheritance rule must be chosen. It is exposed as a
parameter, \code{age\_rule}, precisely because the isochrone analysis depends
on it and the choice must be reported and tested:
\[
  t_{\rm born} \;=\;
  \begin{cases}
    \min(t_1,t_2)                    & \text{\code{'min'}: the oldest ancestor wins,}\\[2pt]
    t_1 \ \text{if}\ m_1 \ge m_2     & \text{\code{'heavier'},}\\[2pt]
    \dfrac{m_1 t_1 + m_2 t_2}{m_1+m_2} & \text{\code{'mass\_weighted'}.}
  \end{cases}
\]

\subsection{Locality of the fragmentation rule --- a physics trap}
\label{sec:trick8}

If the rule is ``the heavier of the pair breaks, whatever hit it'', the
disruption rate of a body of mass $m_0$ against a power-law background
$F(m) \propto m^{\alpha}$ is
\begin{equation}
  \nu(m_0) \;\sim\; \int_{\msink}^{m_0} F(m')\,K(m_0,m')\,dm'
           \;\sim\; \int_{\msink/m_0} x^{\alpha}\,dx ,
  \qquad x = m'/m_0 ,
\end{equation}
which \emph{diverges at the lower limit whenever $\alpha < -1$}. The drift is
then set by the smallest particles in the box --- by the sink scale, not by
$m_0$. The mean-field closure is void, the drift degenerates to
$dm_0/d\tau \sim -C m_0$ (i.e.\ $\beta = 1$, exponential decay), and the
measured index locks onto $-2$ regardless of the kernel: a spurious
universality that is really a boundary effect.

The cure is the one Dohnanyi (1969) built in --- a catastrophic-disruption
threshold. A grain of dust does not shatter a boulder. Requiring
\begin{equation}
  m_{\rm small} \;\ge\; f\,m_{\rm large},
  \qquad f = O(0.1\text{--}1),
\end{equation}
restricts the integral to $x \in [f,1]$, which is scale free, restores
locality, and recovers $\alpha = -(3+\lambda)/2$.

A \emph{closed} system is immune: its self-similar packet is narrow, so all
partners are already of order $m_0$ and the rule is local by construction.
This is why the closed runs agree with the theory at $f=0$ and the open ones
do not.

\subsection{No resampling --- and what it costs}
\label{sec:num_cost}

What is bought:
\begin{itemize}
  \item $\sum w$ and $\sum wm$ conserved to machine precision; the reported
        \code{mass\_drift} is identically zero and any nonzero value is a
        genuine bug;
  \item no clock subtlety at all, $\Delta t = 2V/R$ with $w=1$;
  \item nothing can bias the result, because nothing is resampled.
\end{itemize}

What is paid, and it is the binding constraint on every run:
\begin{itemize}
  \item \textbf{Coagulation.} In a closed box $N = N_i m_i/m_0$, so keeping
        $\sim100$ particles alive at the top requires $m_0 \lesssim N_i/100$:
        $N_i = 10^6$ buys four decades, $N_i = 10^8$ buys six (and
        $\sim3$~GB of arrays). Pushed past that, a run ends with a handful of
        particles; at $N_i = 2\times10^5$ one test finished with a single
        particle alive and the usable plateau shrank from $5.6$ decades to
        $2.7$.
  \item \textbf{Fragmentation.} Worse, in the other direction: $N$ grows as
        $m_i/m_0$, so a completed cascade holds $N_0\,\minj/\msink$
        particles. Three decades multiply the population a thousandfold, and
        \emph{memory, not physics, ends the run}.
\end{itemize}

Use this variant when the priority is an unimpeachable baseline. Use the
weighted variant, which adds multiplication and down-sampling, when the
priority is range: ten decades there cost $4.5\times10^5$ events and returned
$\alpha = -1.005$ against a theoretical $-1.000$. The two agree wherever both
can run, which is the point of keeping both.

\subsection{Snapshot placement}
\label{sec:trick10}

The closed-system observable is a Riemann sum
$F(m) = \sum_k (dN/dm)(m,t_k)\,\Delta t_k$, and the sampling of $t$ has two
separate jobs that must not be confused. \textbf{Placement} decides what is
\emph{resolved}: for a given $m$ the integrand peaks at $m_0 \sim m$, so
every decade of $m$ needs snapshots at the matching decade of $m_0$.
\textbf{Weight} decides what is \emph{computed}: it must be the actual
$\Delta t$ between consecutive snapshots, whatever the placement.

\begin{center}
\small
\begin{tabular}{@{}l p{7.2cm} l@{}}
\toprule
placement & what happens & measured $\alpha$ (theory $-1$) \\
\midrule
\code{events} & events per unit $m_0$ fall as $m_0^{-2}$, so $\Delta t$
  explodes; $94\%$ of the total weight lands in one snapshot
  & $-1.4 \ldots -2.6$ \\
\code{t\_phys}, narrow & fine over $\lesssim 2$ decades & $-0.97 \ldots -1.02$ \\
\code{t\_phys}, wide & $t \propto m_0^{1-\lambda}$, so uniform $t$ is nearly
  uniform $m_0$; over ten decades the lower five get no sample at all
  & $-0.04 \pm 0.16$ \\
\code{log\_m0} & $k$ samples per decade of $m_0$, everywhere
  & $\mathbf{-1.005}$, $R^2 = 0.9995$ \\
\bottomrule
\end{tabular}
\end{center}

The failure of \code{events} deserves emphasis, because it is silent: with
that placement the plateau finder of Sec.~\ref{sec:num_measure} still reports
a clean one-decade plateau at $-1.375 \pm 0.14$. An internally consistent,
entirely wrong measurement. \emph{Fitting the plateau cannot rescue a broken
estimator}, because it tests whether the curve you computed has a scaling
region, not whether that curve is the right quantity. The independent guard
is the printed weight of the last snapshot: it must be $\ll 1\%$.

Resampling (Sec.~\ref{sec:num_cost}) and placement cure different diseases.
Resampling buys the \emph{range}; placement buys a \emph{correct
measurement} over that range.

\paragraph{Rule.} Closed system $\Rightarrow$ \code{log\_m0}, stride $=$
snapshots per decade ($20$--$30$). Open system $\Rightarrow$ the steady state
is the \emph{last} snapshot, $\Delta t$ never enters the estimator, and any
cadence works.

\section{The sink as a criterion}
\label{sec:num_sink}

A stopping or gating criterion phrased in events is a statement about
\emph{resources}; one phrased in sink absorptions is a statement about the
\emph{physics}. Only the second survives a change of $N$, of mass range, or
of kernel. The engine therefore counts $n_{\rm out}$, the number of physical
particles that have left through the absorbing boundary, and exposes it
three ways.

\begin{center}
\small
\begin{tabular}{@{}l p{8.2cm}@{}}
\toprule
\code{iso\_start\_sink} (default $10$) &
  begin accumulating isochrones only after $n_{\rm out}$ absorptions \\
\code{stop\_sink\_events} (default \code{None}) &
  stop the run at $n_{\rm out}$ absorptions \\
\code{out["n\_out"]} &
  the counter as a per-snapshot time series \\
\bottomrule
\end{tabular}
\end{center}

The gate says: bin no age until the cascade has demonstrably pushed that
many particles across the \emph{whole} inertial range. Its cost follows from
the mass budget --- the box holds $\simeq N\sqrt{\msink}$ in steady state and
$n_{\rm out}$ absorptions carry off $n_{\rm out}\msink$ --- so
\begin{equation}
  E \;\simeq\; N\sqrt{\msink} \;+\; n_{\rm out}\,\msink .
  \label{eq:num_budget}
\end{equation}
The second term is independent of $N$: the gate does not get more expensive
as the run grows, only filling the box does. Equation~\eqref{eq:num_budget}
is derived at $\lambda = 0$ and overestimates the cost for $\lambda > 0$,
where a steeper kernel drives mass to the sink faster per event.

A closed system has no sink, so the armed default is silently disarmed
there; only an explicitly passed value raises. The distinction between ``the
caller typed this'' and ``this came from the signature'' is carried by an
\code{int} subclass, so the value behaves as a plain integer everywhere else.

\section{The measurement layer}
\label{sec:num_measure}

A run produces arrays; turning them into an exponent is a separate stage
with its own failure modes.

\paragraph{Estimator.} Closed system: the age-integrated superposition
$F(m) = \sum_t (dN/dm)\,\Delta t$. Open system: the steady state \emph{is}
the instantaneous spectrum, so the last snapshot is the estimator and
$\Delta t$ never enters.

\paragraph{Local slope and the plateau.} The honest way to look at a measured
spectrum is
\begin{equation}
  \Gamma(m) \;=\; \frac{d\log F}{d\log m},
\end{equation}
computed by least squares on a sliding window. A genuine inertial range is a
\emph{plateau} in $\Gamma$; the contaminated ends show up as the places where
$\Gamma$ bends. The reported index is the mean of $\Gamma$ over the longest
contiguous run on which it stays flat within a tolerance, together with its
scatter and its width in decades. A plateau narrower than about one decade is
not a power law, however tight its error bar; the width is reported next to
the index for exactly that reason. If no run qualifies, the answer is
\code{nan}, which is correct rather than a failure.

\paragraph{Guard band.} Independently, strip $0.5$~dex from each end of the
interval between the two characteristic masses --- $[\minj,\msink]$ for an
open system, $[m_i, m_0^{\max}]$ for a closed one --- and fit a single power
law there. This band is chosen \emph{before} looking at the data. The two
estimates must agree; if they do not, the cascade has not developed over the
range that was assumed.

Fitting the whole array does neither. On one closed coagulation run with
$\lambda=0$ and theory $\alpha=-1$: whole array $-1.71$, guard band $-1.09$,
plateau $-1.16 \pm 0.15$ over $1.3$ decades.

\paragraph{Compensated display.} Panels show
\begin{equation}
  m^2 \frac{dN}{dm} \;=\; \text{mass per logarithmic mass interval},
\end{equation}
whose slope is $\alpha+2$. All fitting is done on the raw $dN/dm$; the
compensation is applied at draw time only. The model line carries the plateau
index and an amplitude obtained as the least-squares intercept at
\emph{fixed} slope, so no second fit can quietly move the exponent, and the
theory line uses the same anchoring --- the two curves then differ in slope
alone.

\paragraph{Prediction.} With kernel homogeneity $\lambda$,
\begin{equation}
  \beta =
    \begin{cases}
      \lambda & \text{closed (mass-conserving packet, } n \propto m_0^{-2}),\\[2pt]
      \dfrac{1+\lambda}{2} & \text{open (stationary background, } n \propto m_0^{\alpha}),
    \end{cases}
  \qquad
  \alpha = -(1+\beta),
  \qquad
  b = \frac{1}{1-\beta}.
\end{equation}
Coagulation and fragmentation share $\beta$, $b$ and $\alpha$ exactly; only
the sign of the drift flips, so the law is read in $\tau$ going up and in
$\tau_\ast - \tau$ going down.

\begin{center}
\footnotesize
\begin{tabular}{@{}l l c c c@{}}
\toprule
kernel & $K(m_1,m_2)$ & $\lambda$ &
  $\alpha_{\rm closed}$, $b_{\rm closed}$ &
  $\alpha_{\rm open}$, $b_{\rm open}$ \\
\midrule
constant   & $1$                                & $0$   & $-1$,\ \ $1$    & $-3/2$,\ \ $2$ \\
sum        & $m_1^{1/3}+m_2^{1/3}$              & $1/3$ & $-4/3$,\ \ $3/2$& $-5/3$,\ \ $3$ \\
geometric  & $(m_1^{1/3}+m_2^{1/3})^2$          & $2/3$ & $-5/3$,\ \ $3$  & $-11/6$,\ \ $6$ \\
additive   & $m_1+m_2$                          & $1$   & $-2$,\ \ $\infty$ & $-2$,\ \ $\infty$ \\
product    & $m_1 m_2$                          & $2$   & \multicolumn{2}{c}{above gelation: no cascade} \\
\bottomrule
\end{tabular}
\end{center}

\section{Validation}
\label{sec:num_validation}

Constant kernel, $\lambda = 0$. \emph{plateau} is the automatic index,
\emph{guard} the a-priori band, \emph{dec} the plateau width in decades.

\begin{center}
\small
\begin{tabular}{@{}l r r r r r r@{}}
\toprule
case & $b$ & $b_{\rm th}$ & plateau & $\pm$ & dec & guard \\
\midrule
1 closed coagulation      & $0.990$ & $1$ & $-1.077$ & $0.126$ & $2.70$ & $-1.035$ \\
2 closed fragmentation    & $1.001$ & $1$ & $-0.978$ & $0.141$ & $3.00$ & $-1.056$ \\
3 open coagulation        & $2.001$ & $2$ & $-1.475$ & $0.095$ & $3.60$ & $-1.498$ \\
4 open fragmentation, $f=0$   & --- & $2$ & $-2.047$ & $0.183$ & $0.80$ & $-2.354$ \\
4 open fragmentation, $f=0.3$ & --- & $2$ & $-1.550$ & $0.091$ & $3.40$ & $-1.517$ \\
\bottomrule
\end{tabular}
\end{center}

Case~1 also reproduces the exact Smoluchowski solution for the constant
kernel, $m_0(t) = m_i(1+n_0K_1t/2)$, to $\max|{\rm ratio}-1| = 7.8\times10^{-14}$
--- a direct test of Eq.~\eqref{eq:num_clock}. Mass drift is identically zero
in all four cases.

Two rows are meant to be read as diagnoses rather than measurements. Case~4
at $f=0$ returns $-2.05$ over $0.80$ decades: the index has locked onto $-2$
exactly as Sec.~\ref{sec:trick8} predicts, and the narrow plateau announces
that it is not a measurement. And the growth exponent $b$ is absent for
case~4 because, at this budget, over $99\%$ of live particles still descend
from the initial condition and inherit $t_{\rm born}=0$; the age axis is then
degenerate --- age equals wall-clock time for everyone at once --- and the
isochrones degrade into ``spectrum versus time''. The same fact shows up as
$M_{\rm out}/M_{\rm in} > 1$: the sink is draining the initial reservoir,
whose mass never entered $M_{\rm in}$.

\section{Practical failure modes}
\label{sec:num_pitfalls}

\paragraph{A mass outside the grid is a hard crash.}
The majorant is built from bin \emph{upper corners}, and $b(m)$ clamps
anything above the top edge into the last bin --- whose corner is then
smaller than the particle itself. Hence $K(m,m) > \fmaj_{B-1,B-1}$ and the
first pair raises \code{Majorant violated}. With a constant kernel this can
never happen, since $K\equiv1$ is flat; for any $\lambda>0$ it is immediate.
Check that the grid brackets $\minj$ and $\msink$ before starting, not after.

\paragraph{An unbounded event budget is safe in exactly one case.}
Removing the event cap is only safe when something else can terminate the
run:

\begin{center}
\small
\begin{tabular}{@{}l l p{7.4cm}@{}}
\toprule
case & unbounded & what stops it, or does not \\
\midrule
closed coagulation   & safe   & \code{stop\_max\_mass}; and $N$ decreases monotonically,
                                so \code{live}~$<2$ is a guaranteed terminator \\
closed fragmentation & unsafe & $N$ grows; the stall guard fires only once
                                \emph{everything} sits below the passive floor,
                                i.e.\ at $N_0 m_i/m_{\rm floor}$ particles \\
open coagulation     & hangs  & $N$ sits on a plateau, so no other condition
                                can ever fire \\
open fragmentation   & unsafe & $N$ grows without bound \\
\bottomrule
\end{tabular}
\end{center}

In the open cases the correct brake is \code{stop\_sink\_events}
(Sec.~\ref{sec:num_sink}); keeping a resource ceiling in addition costs
nothing and bounds a mistuned injection rate. The recorded stop reason then
says which one bit.

\paragraph{Changing the kernel invalidates four constants.}
The homogeneity $\lambda$ and every predicted exponent follow automatically,
and the plateau finder, the guard band and the compensation are all
kernel-agnostic. Four things are not:

\begin{enumerate}
  \item \textbf{Injection rate.} Number balance in steady state is
        $q = \langle K\rangle N^2/(2V)$. Writing $q = N^2/2$ is that formula
        at $\langle K\rangle = 1$. Since $N$ is dominated by the injection
        scale when $\alpha<-1$, most pairs are $(\minj,\minj)$, so
        $\langle K \rangle \sim K(\minj,\minj)$, with heavier partners
        pushing it up by roughly a factor two.
  \item \textbf{Collision time.} $t_c = 1/(Kn)$, not $1/n$.
  \item \textbf{Age binning.} One age bin of width $\Delta_\tau$ dex spans
        $b\,\Delta_\tau$ dex of \emph{mass}. Fix the mass resolution and let
        the age step follow, $\Delta_\tau = 0.30/b$; otherwise at $b=6$ a
        step tuned for $b=2$ jumps almost a decade of mass per bin and a
        two-decade fitting window catches two bins.
  \item \textbf{Exact solutions.} The closed-form Smoluchowski reference and
        the ``$m_0(t)$ must be linear'' check are properties of the constant
        kernel alone.
\end{enumerate}

\paragraph{The $-2$ fingerprint has three distinct causes.}
An index of $-2$ is returned by (i) a genuine $\lambda = 1$ kernel, where it
is the correct answer; (ii) the non-local disruption rule of
Sec.~\ref{sec:trick8}, where it is a boundary effect; and (iii) a clock
advanced per event rather than per trial, Sec.~\ref{sec:trick5}, where it is
a numerical artefact. All three drive $\beta \to 1$ and look identical on the
page. Seeing $-2$ therefore says nothing on its own. The discriminating test
is whether changing the kernel changes the index: if it does not, this is not
universality.

\paragraph{A flat shelf is not an inertial range.}
In closed fragmentation the passive floor freezes everything below it: those
particles never fragment again and pile up in a flat shelf that is not part
of the cascade. Left in the array that shelf is the longest flat run there
is --- $4.3$ decades at $\alpha \simeq 0$ --- and the plateau finder duly
reports it. Restrict the search to $m \ge m_{\rm floor}$.

% ---------------------------------------------- delete when including
\end{document}

%      next to the existing \include{Parts/unified-cascade_v4}.
%      main.tex already loads amsmath, amssymb, graphicx and hyperref; you
%      must additionally add to its preamble:
%          \usepackage{algorithm}
%          \usepackage{algpseudocode}
%          \usepackage{booktabs}
% =====================================================================

% ---------------------------------------------- PREAMBLE -- CUT FROM HERE
\documentclass[a4paper,12pt]{article}

\usepackage[T1]{fontenc}
\usepackage[utf8]{inputenc}
\usepackage[english]{babel}
\usepackage{amsmath,amssymb}
\usepackage{geometry}
\geometry{margin=2.5cm}
\usepackage{graphicx}
\usepackage{booktabs}
\usepackage{algorithm}
\usepackage{algpseudocode}
\usepackage[colorlinks=true,linkcolor=blue,citecolor=blue,urlcolor=blue]{hyperref}

\newcommand{\minj}{m_{\mathrm{inj}}}
\newcommand{\msink}{m_{\mathrm{sink}}}
\newcommand{\fmaj}{f^{\mathrm{maj}}}
\newcommand{\code}[1]{\texttt{#1}}

\begin{document}

\title{The numerical scheme: majorant-frequency Monte Carlo\\ for the mass cascade}
\author{}
\date{}
\maketitle
% ---------------------------------------------- PREAMBLE -- CUT TO HERE

\section{Scope and design principle}
\label{sec:num_scope}

The engine integrates the stochastic mass cascade in four configurations,
\[
  \text{process} \in \{\text{coagulation},\ \text{fragmentation}\},
  \qquad
  \text{system} \in \{\text{closed},\ \text{open}\},
\]
through a \emph{single} code path. Only two things differ between the four
cases:

\begin{enumerate}
  \item what an accepted event does to the two selected particles, and
  \item the boundary conditions --- injection plus an absorbing sink, or
        neither.
\end{enumerate}

Everything else --- majorant construction, pair selection, the clock, the
bookkeeping --- is shared. This is a deliberate constraint, not a
convenience: a discrepancy between two of the four cases cannot then be
blamed on two different integrators.

The variant documented here carries no particle weights. One simulated
particle is one physical particle, $w \equiv 1$ for the whole run, so
$\sum_i w_i$ and $\sum_i w_i m_i$ are conserved to machine precision at
every step. The measured mass drift is identically zero, and any nonzero
value is a bug rather than a statistical excursion. Section~\ref{sec:num_cost}
states what that costs.

\subsection{State}

The state is four parallel arrays over live particles, plus a binning of
mass space:

\begin{center}
\small
\begin{tabular}{@{}ll@{}}
\toprule
\code{mass[i]}       & particle mass $m_i$ \\
\code{inj\_time[i]}  & birth time, used for the isochrones (Sec.~\ref{sec:trick7}) \\
\code{bin\_of[i]}    & index of the mass bin holding $i$ \\
\code{pos\_in\_bin[i]} & position of $i$ inside that bin's index list \\
\midrule
\code{edges[0..B]}   & $B+1$ logarithmically spaced mass-bin edges \\
\code{bins[b]}       & list of particle indices in bin $b$ \\
\code{N\_bin[b]}     & occupancy $N_b$ \\
\bottomrule
\end{tabular}
\end{center}

A particle of mass $m$ belongs to bin
$b(m) = \max\bigl(0,\min(B-1,\ \mathrm{searchsorted}(\text{edges},m)-1)\bigr)$.
The clamping in that expression is a trap and is discussed in
Sec.~\ref{sec:num_pitfalls}.

\section{The baseline algorithm}
\label{sec:num_algorithm}

\subsection{Setup}

\begin{algorithm}[h]
\caption{Setup: bins, majorant table, rate vector}
\label{alg:setup}
\begin{algorithmic}[1]
\State bin every initial particle; form the occupancies $N_b$
\State \textbf{majorant table:} for all $b,c$ \quad
       $\fmaj_{bc} \gets K(\text{edges}[b+1],\ \text{edges}[c+1])$
       \Comment{upper corners}
\State $T_b \gets \sum_c \fmaj_{bc} N_c$
\State $\text{row}_b \gets N_b\,(T_b - \fmaj_{bb})$
\State $R \gets \sum_b \text{row}_b$
\end{algorithmic}
\end{algorithm}

The corner rule $\fmaj_{bc}=K(\text{edges}[b+1],\text{edges}[c+1])$ bounds
$K$ over the whole bin pair \emph{provided $K$ is non-decreasing in both
arguments}. That is the single structural assumption on the kernel, and it
is verified at run time rather than assumed: every accepted trial checks
$K(m_i,m_j) \le \fmaj_{bc}$ and raises if it fails.

By construction
\begin{equation}
  R \;=\; \sum_b N_b\bigl(T_b - \fmaj_{bb}\bigr)
    \;=\; \sum_{i \neq j} \fmaj_{ij}
    \;=\; 2 \sum_{i<j} \fmaj_{ij},
  \label{eq:num_R}
\end{equation}
i.e.\ $R$ runs over \emph{ordered} pairs. The factor two in
Eq.~\eqref{eq:num_R} is the origin of the factor two in the clock,
Eq.~\eqref{eq:num_clock}.

\subsection{Main loop}

\begin{algorithm}[h]
\caption{Main loop: one majorant trial}
\label{alg:mainloop}
\begin{algorithmic}[1]
\While{no stopping condition is met}
  \State $\Delta t \gets 2V/(wR)$; \quad $t \gets t + \Delta t$
         \Comment{the clock, Sec.~\ref{sec:trick5}}
  \State \textbf{injection (open systems):}
         $\nu \gets \nu + q\,\Delta t / w$; \quad
         $k \gets \lfloor \nu \rfloor$; \quad $\nu \gets \nu - k$;
         add $k$ particles of mass $\minj$
  \State draw a row bin $b$ with $P(b) \propto \text{row}_b$
  \State draw a column bin $c$ with
         $P(c) \propto N_c \fmaj_{bc}$, using $(N_b-1)\fmaj_{bb}$ for $c=b$
  \State draw $i$ uniformly from \code{bins[b]}, $j \neq i$ uniformly from \code{bins[c]}
  \State \textbf{thin:} accept with probability $K(m_i,m_j)/\fmaj_{bc}$;
         otherwise \textbf{continue} \Comment{null event}
  \State apply the event (Alg.~\ref{alg:event}); apply the absorbing sink
  \State record a snapshot if the cadence says so (Sec.~\ref{sec:trick10})
\EndWhile
\end{algorithmic}
\end{algorithm}

A rejected trial is not an error: it is the thinning that makes the
dominating process exact. The acceptance ratio, reported as \code{acc} in
the run log, is $O(1)$ by construction of the bin-pair table --- measured
$0.95$ for the geometric kernel on bins of $0.1$~dex, where the corner
overshoot is $10^{0.1\lambda}$ per argument.

\subsection{The event}

\begin{algorithm}[h]
\caption{The accepted event --- the only place the four cases differ}
\label{alg:event}
\begin{algorithmic}[1]
\If{coagulation}
  \State $m_{\mathrm{new}} \gets m_i + m_j$
  \State inherit the age of $i$ from $(i,j)$ by \code{age\_rule} (Sec.~\ref{sec:trick7})
  \State $m_i \gets m_{\mathrm{new}}$; rebin $i$; \textbf{delete} $j$
  \State \textbf{if} open \textbf{and} $m_{\mathrm{new}} \ge \msink$ \textbf{then}
         delete $i$; $M_{\mathrm{out}} \mathrel{+}= w\,m_{\mathrm{new}}$;
         $n_{\mathrm{out}} \mathrel{+}= 1$
\Else \Comment{fragmentation}
  \State $p \gets \arg\max(m_i,m_j)$; \quad $s \gets \arg\min(m_i,m_j)$
  \State \textbf{if} $m_s < f\,m_p$ \textbf{then continue}
         \Comment{disruption threshold, Sec.~\ref{sec:trick8}}
  \State $\xi \sim U(0,1)$; \quad $m_a \gets \xi m_p$, \ $m_b \gets (1-\xi)m_p$
  \State both fragments inherit $t_{\mathrm{born}}$ of the single parent $p$
  \State $m_p \gets m_a$; rebin $p$; \textbf{append} a particle of mass $m_b$
  \State \textbf{if} open: delete every fragment with $m \le \msink$,
         highest index first
\EndIf
\end{algorithmic}
\end{algorithm}

Note the sign asymmetry, which governs everything in
Sec.~\ref{sec:num_pitfalls}: coagulation is $2 \to 1$ and therefore a
\emph{sink} of particle number, fragmentation is $1 \to 2$ and therefore a
\emph{source}.

\section{The ten tricks}
\label{sec:num_tricks}

\subsection{Majorant collision frequency}
\label{sec:trick1}

Computing the exact pair rate for every pair is $O(N^2)$ per step. Instead a
majorant $\fmaj \ge K$ drives a dominating Poisson process and events are
thinned with acceptance $K/\fmaj$. The thinned process is exact: the
majorant affects the cost, never the statistics.

\emph{Sources.} Ivanov \& Rogasinsky, Sov.\ J.\ Numer.\ Anal.\ Math.\
Modelling \textbf{3}, 453 (1988) [MCF scheme in DSMC]; Garcia, van den
Broeck, Aertsens \& Serneels, Physica~A \textbf{143}, 535 (1987) [majorant
for coagulation]; Gillespie, J.~Atmos.\ Sci.\ \textbf{32}, 1977 (1975)
[exact stochastic coalescence].

\subsection{Bin-pair majorant table}
\label{sec:trick2}

The majorant is a matrix $\fmaj_{bc}$ over pairs of mass bins, not a single
global number. What that buys is \emph{cost} and \emph{checkability}. It is
worth being precise about which, because a majorant cannot buy correctness:
it cancels out of the physical rate exactly.

\paragraph{The majorant cancels.}
Ticks of the dominating process arrive at rate $wR/(2V)$; a tick selects the
ordered pair $(i,j)$ with probability $\fmaj_{ij}/R$; it is accepted with
probability $K_{ij}/\fmaj_{ij}$. The rate of accepted events on that pair is
therefore
\begin{equation}
  \frac{wR}{2V}\cdot\frac{\fmaj_{ij}}{R}\cdot\frac{K_{ij}}{\fmaj_{ij}}
  \;=\; \frac{w\,K_{ij}}{2V},
  \label{eq:num_cancel}
\end{equation}
and summing over unordered pairs gives $w\sum_{i<j}K_{ij}/V$ --- the
Smoluchowski rate, independent of $\fmaj$ altogether. The majorant cancels
twice, once in the selection weight and once in the acceptance. \emph{Any}
valid majorant therefore returns the same physics; a bad one only wastes
trials. Equation~\eqref{eq:num_cancel} holds because the clock is advanced on
every \emph{trial} rather than on every event --- see
Sec.~\ref{sec:trick5}, which spells out the one edit that breaks it.

\paragraph{What the table actually buys.}
A single global $f_{\max} = K(m_{\max},m_{\max})$ is a legitimate majorant and
gives identical statistics. It is nevertheless the wrong choice, for three
reasons:

\begin{itemize}
  \item \textbf{Acceptance.} A typical pair is accepted with probability
        $\langle K\rangle/f_{\max} \simeq (m_0/m_{\max})^{\lambda}$. Across
        $D$ decades that is $10^{-\lambda D}$: for the geometric kernel over
        four decades, $2\times10^{-3}$, i.e.\ some five hundred trials per
        event against $1.05$ with the table. The table instead keeps the
        acceptance $O(1)$ in \emph{every} bin pair, because the bound is
        evaluated at the corners of the pair actually drawn; the overshoot is
        then only $10^{0.1\lambda}$ per argument at $0.1$~dex bins.
  \item \textbf{The stall guard misfires.} \code{max\_stall\_tries} counts
        trials without an accepted event. At an acceptance of
        $2\times10^{-3}$ its default of $2\times10^{6}$ is only
        $\sim\!4000$ events' worth of honest work, so a healthy run can be
        killed with the diagnosis ``acceptance collapsed''.
  \item \textbf{A lagging maximum is a genuine bias.} A \emph{running} global
        maximum that is not refreshed before some particle grows past it
        gives $\fmaj < K$, and the thinning is then silently wrong --- this,
        unlike the two points above, is a correctness failure. The corner
        table is static, and every trial checks $K \le \fmaj$ explicitly, so
        the same situation raises instead of biasing.
\end{itemize}

\emph{Source.} Eibeck \& Wagner, SIAM J.~Sci.\ Comput.\ \textbf{22}, 802
(2000),
\url{https://doi.org/10.1137/S1064827599353488}.

\subsection{$O(1)$ bin membership}
\label{sec:trick3}

\code{bins[b]} is a list of particle indices and \code{pos\_in\_bin[i]} is
the position of $i$ inside it. Removal swaps the last element into the hole,
so insertion and deletion are $O(1)$ with no reallocation. The same
swap-and-pop idiom is used on the particle arrays themselves when a particle
is deleted, which is why every index held across a deletion must be
revalidated.

\subsection{Incremental rate maintenance}
\label{sec:trick4}

When one bin count changes by $\delta$,
\[
  T \mathrel{+}= \delta\, \fmaj_{\cdot b},
  \qquad
  \text{row} \gets N \odot (T - \mathrm{diag}\,\fmaj),
  \qquad
  R \gets \textstyle\sum_b \text{row}_b ,
\]
an $O(B)$ vector operation instead of recomputing the full matrix--vector
product. Selection uses a vectorised \code{cumsum} plus \code{searchsorted}
rather than a Python loop; that loop is otherwise the dominant cost. For
very large $B$ a Fenwick tree gives $O(\log B)$ (Goodson \& Kraft,
J.~Comput.\ Phys.\ \textbf{183}, 210 (2002)).

\subsection{The clock}
\label{sec:trick5}

The step is
\begin{equation}
  \Delta t \;=\; \frac{2V}{w\,R},
  \label{eq:num_clock}
\end{equation}
and three separate things about it are easy to get wrong.

\paragraph{Ordered versus unordered pairs.}
$R$ counts every unordered pair twice, Eq.~\eqref{eq:num_R}, while the
physical majorant rate runs over unordered pairs. Hence the numerator is
$2V$, not $V$. Pair \emph{selection} is unaffected --- both orderings lead to
the same unordered pair and the factors cancel --- so the error is invisible
in the dynamics and rescales $t$ by exactly two, breaking every comparison
with an analytic solution.

\paragraph{Weights.}
Here $w$ is a constant, but it is still written into the step so that the
formula matches the weighted variant exactly and a constant $w \neq 1$ simply
rescales the density. Nothing in this variant ever changes $w$, so the clock
cannot be corrupted by weight bookkeeping. That is the entire reason the
baseline exists.

\paragraph{Quadratic scaling.}
$R \propto N^2$ because the process is binary, so
\begin{equation}
  t \;\simeq\; \frac{2E}{N^2}
  \label{eq:num_t_of_E}
\end{equation}
after $E$ trials. A step linear in $N$ describes a unary process and gives
the wrong $N(t)$. Equation~\eqref{eq:num_t_of_E} has a practical corollary
that is easy to meet by accident: \emph{raising $N$ at fixed event budget
makes the run shorter in physical time, not longer}. Any gate expressed as an
absolute time and tuned at one $N$ is never crossed at $10N$.

\paragraph{Time advances on every trial, not on every event.}
\code{t\_phys} is incremented immediately after $\Delta t$ is formed, some
fifty lines before the acceptance test, so null (thinned) trials move the
clock too. That is precisely what makes the cancellation
\eqref{eq:num_cancel} work, and moving the increment inside the accepted
branch is the one edit that turns the majorant from a cost into a
\emph{bias}. Physical time would then advance by $2V/(wR)$ per accepted
event, so the reported event rate would be $R$ instead of
$R\,\langle K/\fmaj\rangle = \sum K$, and with a global $f_{\max}$,
\[
  \Delta t_{\rm eff} \;\propto\; \frac{1}{N^2 f_{\max}}
                     \;\propto\; \frac{1}{N^2 m_{\max}^{\lambda}} .
\]
The kernel would then enter through the \emph{extreme} of the distribution
instead of through the typical mass. In an open system $m_{\max}$ is pinned
at the sink and does not grow, so the $m_0$ dependence leaves the clock
entirely, the drift degenerates to $dm_0/dt \propto m_0$ (exponential,
$\beta = 1$), and the constant-flux argument then returns $\alpha = -2$ for
every kernel. This is a third, purely numerical route to the same $-2$
fingerprint that Sec.~\ref{sec:trick8} produces physically.

\subsection{Deterministic residual injection}
\label{sec:trick6}

Injection uses a fractional accumulator,
\[
  \nu \mathrel{+}= q\,\Delta t/w,
  \qquad k = \lfloor \nu \rfloor,
  \qquad \nu \mathrel{-}= k,
\]
rather than a Poisson draw. This removes shot noise at the injection scale.
It is a modelling choice, and it is recorded in the output metadata rather
than left implicit.

\subsection{Isochrone bookkeeping}
\label{sec:trick7}

Every particle carries its birth time; its age is $\tau = t - t_{\rm born}$.
Snapshots accumulate a two-dimensional histogram in $(\tau, m)$, whose rows
are the isochrones and whose sum over rows is the time-averaged spectrum.

For \textbf{fragmentation} the age of a fragment is unambiguous: fragments
descend from exactly one parent. For \textbf{coagulation} a merged particle
has two ancestors and an inheritance rule must be chosen. It is exposed as a
parameter, \code{age\_rule}, precisely because the isochrone analysis depends
on it and the choice must be reported and tested:
\[
  t_{\rm born} \;=\;
  \begin{cases}
    \min(t_1,t_2)                    & \text{\code{'min'}: the oldest ancestor wins,}\\[2pt]
    t_1 \ \text{if}\ m_1 \ge m_2     & \text{\code{'heavier'},}\\[2pt]
    \dfrac{m_1 t_1 + m_2 t_2}{m_1+m_2} & \text{\code{'mass\_weighted'}.}
  \end{cases}
\]

\subsection{Locality of the fragmentation rule --- a physics trap}
\label{sec:trick8}

If the rule is ``the heavier of the pair breaks, whatever hit it'', the
disruption rate of a body of mass $m_0$ against a power-law background
$F(m) \propto m^{\alpha}$ is
\begin{equation}
  \nu(m_0) \;\sim\; \int_{\msink}^{m_0} F(m')\,K(m_0,m')\,dm'
           \;\sim\; \int_{\msink/m_0} x^{\alpha}\,dx ,
  \qquad x = m'/m_0 ,
\end{equation}
which \emph{diverges at the lower limit whenever $\alpha < -1$}. The drift is
then set by the smallest particles in the box --- by the sink scale, not by
$m_0$. The mean-field closure is void, the drift degenerates to
$dm_0/d\tau \sim -C m_0$ (i.e.\ $\beta = 1$, exponential decay), and the
measured index locks onto $-2$ regardless of the kernel: a spurious
universality that is really a boundary effect.

The cure is the one Dohnanyi (1969) built in --- a catastrophic-disruption
threshold. A grain of dust does not shatter a boulder. Requiring
\begin{equation}
  m_{\rm small} \;\ge\; f\,m_{\rm large},
  \qquad f = O(0.1\text{--}1),
\end{equation}
restricts the integral to $x \in [f,1]$, which is scale free, restores
locality, and recovers $\alpha = -(3+\lambda)/2$.

A \emph{closed} system is immune: its self-similar packet is narrow, so all
partners are already of order $m_0$ and the rule is local by construction.
This is why the closed runs agree with the theory at $f=0$ and the open ones
do not.

\subsection{No resampling --- and what it costs}
\label{sec:num_cost}

What is bought:
\begin{itemize}
  \item $\sum w$ and $\sum wm$ conserved to machine precision; the reported
        \code{mass\_drift} is identically zero and any nonzero value is a
        genuine bug;
  \item no clock subtlety at all, $\Delta t = 2V/R$ with $w=1$;
  \item nothing can bias the result, because nothing is resampled.
\end{itemize}

What is paid, and it is the binding constraint on every run:
\begin{itemize}
  \item \textbf{Coagulation.} In a closed box $N = N_i m_i/m_0$, so keeping
        $\sim100$ particles alive at the top requires $m_0 \lesssim N_i/100$:
        $N_i = 10^6$ buys four decades, $N_i = 10^8$ buys six (and
        $\sim3$~GB of arrays). Pushed past that, a run ends with a handful of
        particles; at $N_i = 2\times10^5$ one test finished with a single
        particle alive and the usable plateau shrank from $5.6$ decades to
        $2.7$.
  \item \textbf{Fragmentation.} Worse, in the other direction: $N$ grows as
        $m_i/m_0$, so a completed cascade holds $N_0\,\minj/\msink$
        particles. Three decades multiply the population a thousandfold, and
        \emph{memory, not physics, ends the run}.
\end{itemize}

Use this variant when the priority is an unimpeachable baseline. Use the
weighted variant, which adds multiplication and down-sampling, when the
priority is range: ten decades there cost $4.5\times10^5$ events and returned
$\alpha = -1.005$ against a theoretical $-1.000$. The two agree wherever both
can run, which is the point of keeping both.

\subsection{Snapshot placement}
\label{sec:trick10}

The closed-system observable is a Riemann sum
$F(m) = \sum_k (dN/dm)(m,t_k)\,\Delta t_k$, and the sampling of $t$ has two
separate jobs that must not be confused. \textbf{Placement} decides what is
\emph{resolved}: for a given $m$ the integrand peaks at $m_0 \sim m$, so
every decade of $m$ needs snapshots at the matching decade of $m_0$.
\textbf{Weight} decides what is \emph{computed}: it must be the actual
$\Delta t$ between consecutive snapshots, whatever the placement.

\begin{center}
\small
\begin{tabular}{@{}l p{7.2cm} l@{}}
\toprule
placement & what happens & measured $\alpha$ (theory $-1$) \\
\midrule
\code{events} & events per unit $m_0$ fall as $m_0^{-2}$, so $\Delta t$
  explodes; $94\%$ of the total weight lands in one snapshot
  & $-1.4 \ldots -2.6$ \\
\code{t\_phys}, narrow & fine over $\lesssim 2$ decades & $-0.97 \ldots -1.02$ \\
\code{t\_phys}, wide & $t \propto m_0^{1-\lambda}$, so uniform $t$ is nearly
  uniform $m_0$; over ten decades the lower five get no sample at all
  & $-0.04 \pm 0.16$ \\
\code{log\_m0} & $k$ samples per decade of $m_0$, everywhere
  & $\mathbf{-1.005}$, $R^2 = 0.9995$ \\
\bottomrule
\end{tabular}
\end{center}

The failure of \code{events} deserves emphasis, because it is silent: with
that placement the plateau finder of Sec.~\ref{sec:num_measure} still reports
a clean one-decade plateau at $-1.375 \pm 0.14$. An internally consistent,
entirely wrong measurement. \emph{Fitting the plateau cannot rescue a broken
estimator}, because it tests whether the curve you computed has a scaling
region, not whether that curve is the right quantity. The independent guard
is the printed weight of the last snapshot: it must be $\ll 1\%$.

Resampling (Sec.~\ref{sec:num_cost}) and placement cure different diseases.
Resampling buys the \emph{range}; placement buys a \emph{correct
measurement} over that range.

\paragraph{Rule.} Closed system $\Rightarrow$ \code{log\_m0}, stride $=$
snapshots per decade ($20$--$30$). Open system $\Rightarrow$ the steady state
is the \emph{last} snapshot, $\Delta t$ never enters the estimator, and any
cadence works.

\section{The sink as a criterion}
\label{sec:num_sink}

A stopping or gating criterion phrased in events is a statement about
\emph{resources}; one phrased in sink absorptions is a statement about the
\emph{physics}. Only the second survives a change of $N$, of mass range, or
of kernel. The engine therefore counts $n_{\rm out}$, the number of physical
particles that have left through the absorbing boundary, and exposes it
three ways.

\begin{center}
\small
\begin{tabular}{@{}l p{8.2cm}@{}}
\toprule
\code{iso\_start\_sink} (default $10$) &
  begin accumulating isochrones only after $n_{\rm out}$ absorptions \\
\code{stop\_sink\_events} (default \code{None}) &
  stop the run at $n_{\rm out}$ absorptions \\
\code{out["n\_out"]} &
  the counter as a per-snapshot time series \\
\bottomrule
\end{tabular}
\end{center}

The gate says: bin no age until the cascade has demonstrably pushed that
many particles across the \emph{whole} inertial range. Its cost follows from
the mass budget --- the box holds $\simeq N\sqrt{\msink}$ in steady state and
$n_{\rm out}$ absorptions carry off $n_{\rm out}\msink$ --- so
\begin{equation}
  E \;\simeq\; N\sqrt{\msink} \;+\; n_{\rm out}\,\msink .
  \label{eq:num_budget}
\end{equation}
The second term is independent of $N$: the gate does not get more expensive
as the run grows, only filling the box does. Equation~\eqref{eq:num_budget}
is derived at $\lambda = 0$ and overestimates the cost for $\lambda > 0$,
where a steeper kernel drives mass to the sink faster per event.

A closed system has no sink, so the armed default is silently disarmed
there; only an explicitly passed value raises. The distinction between ``the
caller typed this'' and ``this came from the signature'' is carried by an
\code{int} subclass, so the value behaves as a plain integer everywhere else.

\section{The measurement layer}
\label{sec:num_measure}

A run produces arrays; turning them into an exponent is a separate stage
with its own failure modes.

\paragraph{Estimator.} Closed system: the age-integrated superposition
$F(m) = \sum_t (dN/dm)\,\Delta t$. Open system: the steady state \emph{is}
the instantaneous spectrum, so the last snapshot is the estimator and
$\Delta t$ never enters.

\paragraph{Local slope and the plateau.} The honest way to look at a measured
spectrum is
\begin{equation}
  \Gamma(m) \;=\; \frac{d\log F}{d\log m},
\end{equation}
computed by least squares on a sliding window. A genuine inertial range is a
\emph{plateau} in $\Gamma$; the contaminated ends show up as the places where
$\Gamma$ bends. The reported index is the mean of $\Gamma$ over the longest
contiguous run on which it stays flat within a tolerance, together with its
scatter and its width in decades. A plateau narrower than about one decade is
not a power law, however tight its error bar; the width is reported next to
the index for exactly that reason. If no run qualifies, the answer is
\code{nan}, which is correct rather than a failure.

\paragraph{Guard band.} Independently, strip $0.5$~dex from each end of the
interval between the two characteristic masses --- $[\minj,\msink]$ for an
open system, $[m_i, m_0^{\max}]$ for a closed one --- and fit a single power
law there. This band is chosen \emph{before} looking at the data. The two
estimates must agree; if they do not, the cascade has not developed over the
range that was assumed.

Fitting the whole array does neither. On one closed coagulation run with
$\lambda=0$ and theory $\alpha=-1$: whole array $-1.71$, guard band $-1.09$,
plateau $-1.16 \pm 0.15$ over $1.3$ decades.

\paragraph{Compensated display.} Panels show
\begin{equation}
  m^2 \frac{dN}{dm} \;=\; \text{mass per logarithmic mass interval},
\end{equation}
whose slope is $\alpha+2$. All fitting is done on the raw $dN/dm$; the
compensation is applied at draw time only. The model line carries the plateau
index and an amplitude obtained as the least-squares intercept at
\emph{fixed} slope, so no second fit can quietly move the exponent, and the
theory line uses the same anchoring --- the two curves then differ in slope
alone.

\paragraph{Prediction.} With kernel homogeneity $\lambda$,
\begin{equation}
  \beta =
    \begin{cases}
      \lambda & \text{closed (mass-conserving packet, } n \propto m_0^{-2}),\\[2pt]
      \dfrac{1+\lambda}{2} & \text{open (stationary background, } n \propto m_0^{\alpha}),
    \end{cases}
  \qquad
  \alpha = -(1+\beta),
  \qquad
  b = \frac{1}{1-\beta}.
\end{equation}
Coagulation and fragmentation share $\beta$, $b$ and $\alpha$ exactly; only
the sign of the drift flips, so the law is read in $\tau$ going up and in
$\tau_\ast - \tau$ going down.

\begin{center}
\footnotesize
\begin{tabular}{@{}l l c c c@{}}
\toprule
kernel & $K(m_1,m_2)$ & $\lambda$ &
  $\alpha_{\rm closed}$, $b_{\rm closed}$ &
  $\alpha_{\rm open}$, $b_{\rm open}$ \\
\midrule
constant   & $1$                                & $0$   & $-1$,\ \ $1$    & $-3/2$,\ \ $2$ \\
sum        & $m_1^{1/3}+m_2^{1/3}$              & $1/3$ & $-4/3$,\ \ $3/2$& $-5/3$,\ \ $3$ \\
geometric  & $(m_1^{1/3}+m_2^{1/3})^2$          & $2/3$ & $-5/3$,\ \ $3$  & $-11/6$,\ \ $6$ \\
additive   & $m_1+m_2$                          & $1$   & $-2$,\ \ $\infty$ & $-2$,\ \ $\infty$ \\
product    & $m_1 m_2$                          & $2$   & \multicolumn{2}{c}{above gelation: no cascade} \\
\bottomrule
\end{tabular}
\end{center}

\section{Validation}
\label{sec:num_validation}

Constant kernel, $\lambda = 0$. \emph{plateau} is the automatic index,
\emph{guard} the a-priori band, \emph{dec} the plateau width in decades.

\begin{center}
\small
\begin{tabular}{@{}l r r r r r r@{}}
\toprule
case & $b$ & $b_{\rm th}$ & plateau & $\pm$ & dec & guard \\
\midrule
1 closed coagulation      & $0.990$ & $1$ & $-1.077$ & $0.126$ & $2.70$ & $-1.035$ \\
2 closed fragmentation    & $1.001$ & $1$ & $-0.978$ & $0.141$ & $3.00$ & $-1.056$ \\
3 open coagulation        & $2.001$ & $2$ & $-1.475$ & $0.095$ & $3.60$ & $-1.498$ \\
4 open fragmentation, $f=0$   & --- & $2$ & $-2.047$ & $0.183$ & $0.80$ & $-2.354$ \\
4 open fragmentation, $f=0.3$ & --- & $2$ & $-1.550$ & $0.091$ & $3.40$ & $-1.517$ \\
\bottomrule
\end{tabular}
\end{center}

Case~1 also reproduces the exact Smoluchowski solution for the constant
kernel, $m_0(t) = m_i(1+n_0K_1t/2)$, to $\max|{\rm ratio}-1| = 7.8\times10^{-14}$
--- a direct test of Eq.~\eqref{eq:num_clock}. Mass drift is identically zero
in all four cases.

Two rows are meant to be read as diagnoses rather than measurements. Case~4
at $f=0$ returns $-2.05$ over $0.80$ decades: the index has locked onto $-2$
exactly as Sec.~\ref{sec:trick8} predicts, and the narrow plateau announces
that it is not a measurement. And the growth exponent $b$ is absent for
case~4 because, at this budget, over $99\%$ of live particles still descend
from the initial condition and inherit $t_{\rm born}=0$; the age axis is then
degenerate --- age equals wall-clock time for everyone at once --- and the
isochrones degrade into ``spectrum versus time''. The same fact shows up as
$M_{\rm out}/M_{\rm in} > 1$: the sink is draining the initial reservoir,
whose mass never entered $M_{\rm in}$.

\section{Practical failure modes}
\label{sec:num_pitfalls}

\paragraph{A mass outside the grid is a hard crash.}
The majorant is built from bin \emph{upper corners}, and $b(m)$ clamps
anything above the top edge into the last bin --- whose corner is then
smaller than the particle itself. Hence $K(m,m) > \fmaj_{B-1,B-1}$ and the
first pair raises \code{Majorant violated}. With a constant kernel this can
never happen, since $K\equiv1$ is flat; for any $\lambda>0$ it is immediate.
Check that the grid brackets $\minj$ and $\msink$ before starting, not after.

\paragraph{An unbounded event budget is safe in exactly one case.}
Removing the event cap is only safe when something else can terminate the
run:

\begin{center}
\small
\begin{tabular}{@{}l l p{7.4cm}@{}}
\toprule
case & unbounded & what stops it, or does not \\
\midrule
closed coagulation   & safe   & \code{stop\_max\_mass}; and $N$ decreases monotonically,
                                so \code{live}~$<2$ is a guaranteed terminator \\
closed fragmentation & unsafe & $N$ grows; the stall guard fires only once
                                \emph{everything} sits below the passive floor,
                                i.e.\ at $N_0 m_i/m_{\rm floor}$ particles \\
open coagulation     & hangs  & $N$ sits on a plateau, so no other condition
                                can ever fire \\
open fragmentation   & unsafe & $N$ grows without bound \\
\bottomrule
\end{tabular}
\end{center}

In the open cases the correct brake is \code{stop\_sink\_events}
(Sec.~\ref{sec:num_sink}); keeping a resource ceiling in addition costs
nothing and bounds a mistuned injection rate. The recorded stop reason then
says which one bit.

\paragraph{Changing the kernel invalidates four constants.}
The homogeneity $\lambda$ and every predicted exponent follow automatically,
and the plateau finder, the guard band and the compensation are all
kernel-agnostic. Four things are not:

\begin{enumerate}
  \item \textbf{Injection rate.} Number balance in steady state is
        $q = \langle K\rangle N^2/(2V)$. Writing $q = N^2/2$ is that formula
        at $\langle K\rangle = 1$. Since $N$ is dominated by the injection
        scale when $\alpha<-1$, most pairs are $(\minj,\minj)$, so
        $\langle K \rangle \sim K(\minj,\minj)$, with heavier partners
        pushing it up by roughly a factor two.
  \item \textbf{Collision time.} $t_c = 1/(Kn)$, not $1/n$.
  \item \textbf{Age binning.} One age bin of width $\Delta_\tau$ dex spans
        $b\,\Delta_\tau$ dex of \emph{mass}. Fix the mass resolution and let
        the age step follow, $\Delta_\tau = 0.30/b$; otherwise at $b=6$ a
        step tuned for $b=2$ jumps almost a decade of mass per bin and a
        two-decade fitting window catches two bins.
  \item \textbf{Exact solutions.} The closed-form Smoluchowski reference and
        the ``$m_0(t)$ must be linear'' check are properties of the constant
        kernel alone.
\end{enumerate}

\paragraph{The $-2$ fingerprint has three distinct causes.}
An index of $-2$ is returned by (i) a genuine $\lambda = 1$ kernel, where it
is the correct answer; (ii) the non-local disruption rule of
Sec.~\ref{sec:trick8}, where it is a boundary effect; and (iii) a clock
advanced per event rather than per trial, Sec.~\ref{sec:trick5}, where it is
a numerical artefact. All three drive $\beta \to 1$ and look identical on the
page. Seeing $-2$ therefore says nothing on its own. The discriminating test
is whether changing the kernel changes the index: if it does not, this is not
universality.

\paragraph{A flat shelf is not an inertial range.}
In closed fragmentation the passive floor freezes everything below it: those
particles never fragment again and pile up in a flat shelf that is not part
of the cascade. Left in the array that shelf is the longest flat run there
is --- $4.3$ decades at $\alpha \simeq 0$ --- and the plateau finder duly
reports it. Restrict the search to $m \ge m_{\rm floor}$.

% ---------------------------------------------- delete when including
\end{document}

%      next to the existing \include{Parts/unified-cascade_v4}.
%      main.tex already loads amsmath, amssymb, graphicx and hyperref; you
%      must additionally add to its preamble:
%          \usepackage{algorithm}
%          \usepackage{algpseudocode}
%          \usepackage{booktabs}
% =====================================================================

% ---------------------------------------------- PREAMBLE -- CUT FROM HERE
\documentclass[a4paper,12pt]{article}

\usepackage[T1]{fontenc}
\usepackage[utf8]{inputenc}
\usepackage[english]{babel}
\usepackage{amsmath,amssymb}
\usepackage{geometry}
\geometry{margin=2.5cm}
\usepackage{graphicx}
\usepackage{booktabs}
\usepackage{algorithm}
\usepackage{algpseudocode}
\usepackage[colorlinks=true,linkcolor=blue,citecolor=blue,urlcolor=blue]{hyperref}

\newcommand{\minj}{m_{\mathrm{inj}}}
\newcommand{\msink}{m_{\mathrm{sink}}}
\newcommand{\fmaj}{f^{\mathrm{maj}}}
\newcommand{\code}[1]{\texttt{#1}}

\begin{document}

\title{The numerical scheme: majorant-frequency Monte Carlo\\ for the mass cascade}
\author{}
\date{}
\maketitle
% ---------------------------------------------- PREAMBLE -- CUT TO HERE

\section{Scope and design principle}
\label{sec:num_scope}

The engine integrates the stochastic mass cascade in four configurations,
\[
  \text{process} \in \{\text{coagulation},\ \text{fragmentation}\},
  \qquad
  \text{system} \in \{\text{closed},\ \text{open}\},
\]
through a \emph{single} code path. Only two things differ between the four
cases:

\begin{enumerate}
  \item what an accepted event does to the two selected particles, and
  \item the boundary conditions --- injection plus an absorbing sink, or
        neither.
\end{enumerate}

Everything else --- majorant construction, pair selection, the clock, the
bookkeeping --- is shared. This is a deliberate constraint, not a
convenience: a discrepancy between two of the four cases cannot then be
blamed on two different integrators.

The variant documented here carries no particle weights. One simulated
particle is one physical particle, $w \equiv 1$ for the whole run, so
$\sum_i w_i$ and $\sum_i w_i m_i$ are conserved to machine precision at
every step. The measured mass drift is identically zero, and any nonzero
value is a bug rather than a statistical excursion. Section~\ref{sec:num_cost}
states what that costs.

\subsection{State}

The state is four parallel arrays over live particles, plus a binning of
mass space:

\begin{center}
\small
\begin{tabular}{@{}ll@{}}
\toprule
\code{mass[i]}       & particle mass $m_i$ \\
\code{inj\_time[i]}  & birth time, used for the isochrones (Sec.~\ref{sec:trick7}) \\
\code{bin\_of[i]}    & index of the mass bin holding $i$ \\
\code{pos\_in\_bin[i]} & position of $i$ inside that bin's index list \\
\midrule
\code{edges[0..B]}   & $B+1$ logarithmically spaced mass-bin edges \\
\code{bins[b]}       & list of particle indices in bin $b$ \\
\code{N\_bin[b]}     & occupancy $N_b$ \\
\bottomrule
\end{tabular}
\end{center}

A particle of mass $m$ belongs to bin
$b(m) = \max\bigl(0,\min(B-1,\ \mathrm{searchsorted}(\text{edges},m)-1)\bigr)$.
The clamping in that expression is a trap and is discussed in
Sec.~\ref{sec:num_pitfalls}.

\section{The baseline algorithm}
\label{sec:num_algorithm}

\subsection{Setup}

\begin{algorithm}[h]
\caption{Setup: bins, majorant table, rate vector}
\label{alg:setup}
\begin{algorithmic}[1]
\State bin every initial particle; form the occupancies $N_b$
\State \textbf{majorant table:} for all $b,c$ \quad
       $\fmaj_{bc} \gets K(\text{edges}[b+1],\ \text{edges}[c+1])$
       \Comment{upper corners}
\State $T_b \gets \sum_c \fmaj_{bc} N_c$
\State $\text{row}_b \gets N_b\,(T_b - \fmaj_{bb})$
\State $R \gets \sum_b \text{row}_b$
\end{algorithmic}
\end{algorithm}

The corner rule $\fmaj_{bc}=K(\text{edges}[b+1],\text{edges}[c+1])$ bounds
$K$ over the whole bin pair \emph{provided $K$ is non-decreasing in both
arguments}. That is the single structural assumption on the kernel, and it
is verified at run time rather than assumed: every accepted trial checks
$K(m_i,m_j) \le \fmaj_{bc}$ and raises if it fails.

By construction
\begin{equation}
  R \;=\; \sum_b N_b\bigl(T_b - \fmaj_{bb}\bigr)
    \;=\; \sum_{i \neq j} \fmaj_{ij}
    \;=\; 2 \sum_{i<j} \fmaj_{ij},
  \label{eq:num_R}
\end{equation}
i.e.\ $R$ runs over \emph{ordered} pairs. The factor two in
Eq.~\eqref{eq:num_R} is the origin of the factor two in the clock,
Eq.~\eqref{eq:num_clock}.

\subsection{Main loop}

\begin{algorithm}[h]
\caption{Main loop: one majorant trial}
\label{alg:mainloop}
\begin{algorithmic}[1]
\While{no stopping condition is met}
  \State $\Delta t \gets 2V/(wR)$; \quad $t \gets t + \Delta t$
         \Comment{the clock, Sec.~\ref{sec:trick5}}
  \State \textbf{injection (open systems):}
         $\nu \gets \nu + q\,\Delta t / w$; \quad
         $k \gets \lfloor \nu \rfloor$; \quad $\nu \gets \nu - k$;
         add $k$ particles of mass $\minj$
  \State draw a row bin $b$ with $P(b) \propto \text{row}_b$
  \State draw a column bin $c$ with
         $P(c) \propto N_c \fmaj_{bc}$, using $(N_b-1)\fmaj_{bb}$ for $c=b$
  \State draw $i$ uniformly from \code{bins[b]}, $j \neq i$ uniformly from \code{bins[c]}
  \State \textbf{thin:} accept with probability $K(m_i,m_j)/\fmaj_{bc}$;
         otherwise \textbf{continue} \Comment{null event}
  \State apply the event (Alg.~\ref{alg:event}); apply the absorbing sink
  \State record a snapshot if the cadence says so (Sec.~\ref{sec:trick10})
\EndWhile
\end{algorithmic}
\end{algorithm}

A rejected trial is not an error: it is the thinning that makes the
dominating process exact. The acceptance ratio, reported as \code{acc} in
the run log, is $O(1)$ by construction of the bin-pair table --- measured
$0.95$ for the geometric kernel on bins of $0.1$~dex, where the corner
overshoot is $10^{0.1\lambda}$ per argument.

\subsection{The event}

\begin{algorithm}[h]
\caption{The accepted event --- the only place the four cases differ}
\label{alg:event}
\begin{algorithmic}[1]
\If{coagulation}
  \State $m_{\mathrm{new}} \gets m_i + m_j$
  \State inherit the age of $i$ from $(i,j)$ by \code{age\_rule} (Sec.~\ref{sec:trick7})
  \State $m_i \gets m_{\mathrm{new}}$; rebin $i$; \textbf{delete} $j$
  \State \textbf{if} open \textbf{and} $m_{\mathrm{new}} \ge \msink$ \textbf{then}
         delete $i$; $M_{\mathrm{out}} \mathrel{+}= w\,m_{\mathrm{new}}$;
         $n_{\mathrm{out}} \mathrel{+}= 1$
\Else \Comment{fragmentation}
  \State $p \gets \arg\max(m_i,m_j)$; \quad $s \gets \arg\min(m_i,m_j)$
  \State \textbf{if} $m_s < f\,m_p$ \textbf{then continue}
         \Comment{disruption threshold, Sec.~\ref{sec:trick8}}
  \State $\xi \sim U(0,1)$; \quad $m_a \gets \xi m_p$, \ $m_b \gets (1-\xi)m_p$
  \State both fragments inherit $t_{\mathrm{born}}$ of the single parent $p$
  \State $m_p \gets m_a$; rebin $p$; \textbf{append} a particle of mass $m_b$
  \State \textbf{if} open: delete every fragment with $m \le \msink$,
         highest index first
\EndIf
\end{algorithmic}
\end{algorithm}

Note the sign asymmetry, which governs everything in
Sec.~\ref{sec:num_pitfalls}: coagulation is $2 \to 1$ and therefore a
\emph{sink} of particle number, fragmentation is $1 \to 2$ and therefore a
\emph{source}.

\section{The ten tricks}
\label{sec:num_tricks}

\subsection{Majorant collision frequency}
\label{sec:trick1}

Computing the exact pair rate for every pair is $O(N^2)$ per step. Instead a
majorant $\fmaj \ge K$ drives a dominating Poisson process and events are
thinned with acceptance $K/\fmaj$. The thinned process is exact: the
majorant affects the cost, never the statistics.

\emph{Sources.} Ivanov \& Rogasinsky, Sov.\ J.\ Numer.\ Anal.\ Math.\
Modelling \textbf{3}, 453 (1988) [MCF scheme in DSMC]; Garcia, van den
Broeck, Aertsens \& Serneels, Physica~A \textbf{143}, 535 (1987) [majorant
for coagulation]; Gillespie, J.~Atmos.\ Sci.\ \textbf{32}, 1977 (1975)
[exact stochastic coalescence].

\subsection{Bin-pair majorant table}
\label{sec:trick2}

The majorant is a matrix $\fmaj_{bc}$ over pairs of mass bins, not a single
global number. What that buys is \emph{cost} and \emph{checkability}. It is
worth being precise about which, because a majorant cannot buy correctness:
it cancels out of the physical rate exactly.

\paragraph{The majorant cancels.}
Ticks of the dominating process arrive at rate $wR/(2V)$; a tick selects the
ordered pair $(i,j)$ with probability $\fmaj_{ij}/R$; it is accepted with
probability $K_{ij}/\fmaj_{ij}$. The rate of accepted events on that pair is
therefore
\begin{equation}
  \frac{wR}{2V}\cdot\frac{\fmaj_{ij}}{R}\cdot\frac{K_{ij}}{\fmaj_{ij}}
  \;=\; \frac{w\,K_{ij}}{2V},
  \label{eq:num_cancel}
\end{equation}
and summing over unordered pairs gives $w\sum_{i<j}K_{ij}/V$ --- the
Smoluchowski rate, independent of $\fmaj$ altogether. The majorant cancels
twice, once in the selection weight and once in the acceptance. \emph{Any}
valid majorant therefore returns the same physics; a bad one only wastes
trials. Equation~\eqref{eq:num_cancel} holds because the clock is advanced on
every \emph{trial} rather than on every event --- see
Sec.~\ref{sec:trick5}, which spells out the one edit that breaks it.

\paragraph{What the table actually buys.}
A single global $f_{\max} = K(m_{\max},m_{\max})$ is a legitimate majorant and
gives identical statistics. It is nevertheless the wrong choice, for three
reasons:

\begin{itemize}
  \item \textbf{Acceptance.} A typical pair is accepted with probability
        $\langle K\rangle/f_{\max} \simeq (m_0/m_{\max})^{\lambda}$. Across
        $D$ decades that is $10^{-\lambda D}$: for the geometric kernel over
        four decades, $2\times10^{-3}$, i.e.\ some five hundred trials per
        event against $1.05$ with the table. The table instead keeps the
        acceptance $O(1)$ in \emph{every} bin pair, because the bound is
        evaluated at the corners of the pair actually drawn; the overshoot is
        then only $10^{0.1\lambda}$ per argument at $0.1$~dex bins.
  \item \textbf{The stall guard misfires.} \code{max\_stall\_tries} counts
        trials without an accepted event. At an acceptance of
        $2\times10^{-3}$ its default of $2\times10^{6}$ is only
        $\sim\!4000$ events' worth of honest work, so a healthy run can be
        killed with the diagnosis ``acceptance collapsed''.
  \item \textbf{A lagging maximum is a genuine bias.} A \emph{running} global
        maximum that is not refreshed before some particle grows past it
        gives $\fmaj < K$, and the thinning is then silently wrong --- this,
        unlike the two points above, is a correctness failure. The corner
        table is static, and every trial checks $K \le \fmaj$ explicitly, so
        the same situation raises instead of biasing.
\end{itemize}

\emph{Source.} Eibeck \& Wagner, SIAM J.~Sci.\ Comput.\ \textbf{22}, 802
(2000),
\url{https://doi.org/10.1137/S1064827599353488}.

\subsection{$O(1)$ bin membership}
\label{sec:trick3}

\code{bins[b]} is a list of particle indices and \code{pos\_in\_bin[i]} is
the position of $i$ inside it. Removal swaps the last element into the hole,
so insertion and deletion are $O(1)$ with no reallocation. The same
swap-and-pop idiom is used on the particle arrays themselves when a particle
is deleted, which is why every index held across a deletion must be
revalidated.

\subsection{Incremental rate maintenance}
\label{sec:trick4}

When one bin count changes by $\delta$,
\[
  T \mathrel{+}= \delta\, \fmaj_{\cdot b},
  \qquad
  \text{row} \gets N \odot (T - \mathrm{diag}\,\fmaj),
  \qquad
  R \gets \textstyle\sum_b \text{row}_b ,
\]
an $O(B)$ vector operation instead of recomputing the full matrix--vector
product. Selection uses a vectorised \code{cumsum} plus \code{searchsorted}
rather than a Python loop; that loop is otherwise the dominant cost. For
very large $B$ a Fenwick tree gives $O(\log B)$ (Goodson \& Kraft,
J.~Comput.\ Phys.\ \textbf{183}, 210 (2002)).

\subsection{The clock}
\label{sec:trick5}

The step is
\begin{equation}
  \Delta t \;=\; \frac{2V}{w\,R},
  \label{eq:num_clock}
\end{equation}
and three separate things about it are easy to get wrong.

\paragraph{Ordered versus unordered pairs.}
$R$ counts every unordered pair twice, Eq.~\eqref{eq:num_R}, while the
physical majorant rate runs over unordered pairs. Hence the numerator is
$2V$, not $V$. Pair \emph{selection} is unaffected --- both orderings lead to
the same unordered pair and the factors cancel --- so the error is invisible
in the dynamics and rescales $t$ by exactly two, breaking every comparison
with an analytic solution.

\paragraph{Weights.}
Here $w$ is a constant, but it is still written into the step so that the
formula matches the weighted variant exactly and a constant $w \neq 1$ simply
rescales the density. Nothing in this variant ever changes $w$, so the clock
cannot be corrupted by weight bookkeeping. That is the entire reason the
baseline exists.

\paragraph{Quadratic scaling.}
$R \propto N^2$ because the process is binary, so
\begin{equation}
  t \;\simeq\; \frac{2E}{N^2}
  \label{eq:num_t_of_E}
\end{equation}
after $E$ trials. A step linear in $N$ describes a unary process and gives
the wrong $N(t)$. Equation~\eqref{eq:num_t_of_E} has a practical corollary
that is easy to meet by accident: \emph{raising $N$ at fixed event budget
makes the run shorter in physical time, not longer}. Any gate expressed as an
absolute time and tuned at one $N$ is never crossed at $10N$.

\paragraph{Time advances on every trial, not on every event.}
\code{t\_phys} is incremented immediately after $\Delta t$ is formed, some
fifty lines before the acceptance test, so null (thinned) trials move the
clock too. That is precisely what makes the cancellation
\eqref{eq:num_cancel} work, and moving the increment inside the accepted
branch is the one edit that turns the majorant from a cost into a
\emph{bias}. Physical time would then advance by $2V/(wR)$ per accepted
event, so the reported event rate would be $R$ instead of
$R\,\langle K/\fmaj\rangle = \sum K$, and with a global $f_{\max}$,
\[
  \Delta t_{\rm eff} \;\propto\; \frac{1}{N^2 f_{\max}}
                     \;\propto\; \frac{1}{N^2 m_{\max}^{\lambda}} .
\]
The kernel would then enter through the \emph{extreme} of the distribution
instead of through the typical mass. In an open system $m_{\max}$ is pinned
at the sink and does not grow, so the $m_0$ dependence leaves the clock
entirely, the drift degenerates to $dm_0/dt \propto m_0$ (exponential,
$\beta = 1$), and the constant-flux argument then returns $\alpha = -2$ for
every kernel. This is a third, purely numerical route to the same $-2$
fingerprint that Sec.~\ref{sec:trick8} produces physically.

\subsection{Deterministic residual injection}
\label{sec:trick6}

Injection uses a fractional accumulator,
\[
  \nu \mathrel{+}= q\,\Delta t/w,
  \qquad k = \lfloor \nu \rfloor,
  \qquad \nu \mathrel{-}= k,
\]
rather than a Poisson draw. This removes shot noise at the injection scale.
It is a modelling choice, and it is recorded in the output metadata rather
than left implicit.

\subsection{Isochrone bookkeeping}
\label{sec:trick7}

Every particle carries its birth time; its age is $\tau = t - t_{\rm born}$.
Snapshots accumulate a two-dimensional histogram in $(\tau, m)$, whose rows
are the isochrones and whose sum over rows is the time-averaged spectrum.

For \textbf{fragmentation} the age of a fragment is unambiguous: fragments
descend from exactly one parent. For \textbf{coagulation} a merged particle
has two ancestors and an inheritance rule must be chosen. It is exposed as a
parameter, \code{age\_rule}, precisely because the isochrone analysis depends
on it and the choice must be reported and tested:
\[
  t_{\rm born} \;=\;
  \begin{cases}
    \min(t_1,t_2)                    & \text{\code{'min'}: the oldest ancestor wins,}\\[2pt]
    t_1 \ \text{if}\ m_1 \ge m_2     & \text{\code{'heavier'},}\\[2pt]
    \dfrac{m_1 t_1 + m_2 t_2}{m_1+m_2} & \text{\code{'mass\_weighted'}.}
  \end{cases}
\]

\subsection{Locality of the fragmentation rule --- a physics trap}
\label{sec:trick8}

If the rule is ``the heavier of the pair breaks, whatever hit it'', the
disruption rate of a body of mass $m_0$ against a power-law background
$F(m) \propto m^{\alpha}$ is
\begin{equation}
  \nu(m_0) \;\sim\; \int_{\msink}^{m_0} F(m')\,K(m_0,m')\,dm'
           \;\sim\; \int_{\msink/m_0} x^{\alpha}\,dx ,
  \qquad x = m'/m_0 ,
\end{equation}
which \emph{diverges at the lower limit whenever $\alpha < -1$}. The drift is
then set by the smallest particles in the box --- by the sink scale, not by
$m_0$. The mean-field closure is void, the drift degenerates to
$dm_0/d\tau \sim -C m_0$ (i.e.\ $\beta = 1$, exponential decay), and the
measured index locks onto $-2$ regardless of the kernel: a spurious
universality that is really a boundary effect.

The cure is the one Dohnanyi (1969) built in --- a catastrophic-disruption
threshold. A grain of dust does not shatter a boulder. Requiring
\begin{equation}
  m_{\rm small} \;\ge\; f\,m_{\rm large},
  \qquad f = O(0.1\text{--}1),
\end{equation}
restricts the integral to $x \in [f,1]$, which is scale free, restores
locality, and recovers $\alpha = -(3+\lambda)/2$.

A \emph{closed} system is immune: its self-similar packet is narrow, so all
partners are already of order $m_0$ and the rule is local by construction.
This is why the closed runs agree with the theory at $f=0$ and the open ones
do not.

\subsection{No resampling --- and what it costs}
\label{sec:num_cost}

What is bought:
\begin{itemize}
  \item $\sum w$ and $\sum wm$ conserved to machine precision; the reported
        \code{mass\_drift} is identically zero and any nonzero value is a
        genuine bug;
  \item no clock subtlety at all, $\Delta t = 2V/R$ with $w=1$;
  \item nothing can bias the result, because nothing is resampled.
\end{itemize}

What is paid, and it is the binding constraint on every run:
\begin{itemize}
  \item \textbf{Coagulation.} In a closed box $N = N_i m_i/m_0$, so keeping
        $\sim100$ particles alive at the top requires $m_0 \lesssim N_i/100$:
        $N_i = 10^6$ buys four decades, $N_i = 10^8$ buys six (and
        $\sim3$~GB of arrays). Pushed past that, a run ends with a handful of
        particles; at $N_i = 2\times10^5$ one test finished with a single
        particle alive and the usable plateau shrank from $5.6$ decades to
        $2.7$.
  \item \textbf{Fragmentation.} Worse, in the other direction: $N$ grows as
        $m_i/m_0$, so a completed cascade holds $N_0\,\minj/\msink$
        particles. Three decades multiply the population a thousandfold, and
        \emph{memory, not physics, ends the run}.
\end{itemize}

Use this variant when the priority is an unimpeachable baseline. Use the
weighted variant, which adds multiplication and down-sampling, when the
priority is range: ten decades there cost $4.5\times10^5$ events and returned
$\alpha = -1.005$ against a theoretical $-1.000$. The two agree wherever both
can run, which is the point of keeping both.

\subsection{Snapshot placement}
\label{sec:trick10}

The closed-system observable is a Riemann sum
$F(m) = \sum_k (dN/dm)(m,t_k)\,\Delta t_k$, and the sampling of $t$ has two
separate jobs that must not be confused. \textbf{Placement} decides what is
\emph{resolved}: for a given $m$ the integrand peaks at $m_0 \sim m$, so
every decade of $m$ needs snapshots at the matching decade of $m_0$.
\textbf{Weight} decides what is \emph{computed}: it must be the actual
$\Delta t$ between consecutive snapshots, whatever the placement.

\begin{center}
\small
\begin{tabular}{@{}l p{7.2cm} l@{}}
\toprule
placement & what happens & measured $\alpha$ (theory $-1$) \\
\midrule
\code{events} & events per unit $m_0$ fall as $m_0^{-2}$, so $\Delta t$
  explodes; $94\%$ of the total weight lands in one snapshot
  & $-1.4 \ldots -2.6$ \\
\code{t\_phys}, narrow & fine over $\lesssim 2$ decades & $-0.97 \ldots -1.02$ \\
\code{t\_phys}, wide & $t \propto m_0^{1-\lambda}$, so uniform $t$ is nearly
  uniform $m_0$; over ten decades the lower five get no sample at all
  & $-0.04 \pm 0.16$ \\
\code{log\_m0} & $k$ samples per decade of $m_0$, everywhere
  & $\mathbf{-1.005}$, $R^2 = 0.9995$ \\
\bottomrule
\end{tabular}
\end{center}

The failure of \code{events} deserves emphasis, because it is silent: with
that placement the plateau finder of Sec.~\ref{sec:num_measure} still reports
a clean one-decade plateau at $-1.375 \pm 0.14$. An internally consistent,
entirely wrong measurement. \emph{Fitting the plateau cannot rescue a broken
estimator}, because it tests whether the curve you computed has a scaling
region, not whether that curve is the right quantity. The independent guard
is the printed weight of the last snapshot: it must be $\ll 1\%$.

Resampling (Sec.~\ref{sec:num_cost}) and placement cure different diseases.
Resampling buys the \emph{range}; placement buys a \emph{correct
measurement} over that range.

\paragraph{Rule.} Closed system $\Rightarrow$ \code{log\_m0}, stride $=$
snapshots per decade ($20$--$30$). Open system $\Rightarrow$ the steady state
is the \emph{last} snapshot, $\Delta t$ never enters the estimator, and any
cadence works.

\section{The sink as a criterion}
\label{sec:num_sink}

A stopping or gating criterion phrased in events is a statement about
\emph{resources}; one phrased in sink absorptions is a statement about the
\emph{physics}. Only the second survives a change of $N$, of mass range, or
of kernel. The engine therefore counts $n_{\rm out}$, the number of physical
particles that have left through the absorbing boundary, and exposes it
three ways.

\begin{center}
\small
\begin{tabular}{@{}l p{8.2cm}@{}}
\toprule
\code{iso\_start\_sink} (default $10$) &
  begin accumulating isochrones only after $n_{\rm out}$ absorptions \\
\code{stop\_sink\_events} (default \code{None}) &
  stop the run at $n_{\rm out}$ absorptions \\
\code{out["n\_out"]} &
  the counter as a per-snapshot time series \\
\bottomrule
\end{tabular}
\end{center}

The gate says: bin no age until the cascade has demonstrably pushed that
many particles across the \emph{whole} inertial range. Its cost follows from
the mass budget --- the box holds $\simeq N\sqrt{\msink}$ in steady state and
$n_{\rm out}$ absorptions carry off $n_{\rm out}\msink$ --- so
\begin{equation}
  E \;\simeq\; N\sqrt{\msink} \;+\; n_{\rm out}\,\msink .
  \label{eq:num_budget}
\end{equation}
The second term is independent of $N$: the gate does not get more expensive
as the run grows, only filling the box does. Equation~\eqref{eq:num_budget}
is derived at $\lambda = 0$ and overestimates the cost for $\lambda > 0$,
where a steeper kernel drives mass to the sink faster per event.

A closed system has no sink, so the armed default is silently disarmed
there; only an explicitly passed value raises. The distinction between ``the
caller typed this'' and ``this came from the signature'' is carried by an
\code{int} subclass, so the value behaves as a plain integer everywhere else.

\section{The measurement layer}
\label{sec:num_measure}

A run produces arrays; turning them into an exponent is a separate stage
with its own failure modes.

\paragraph{Estimator.} Closed system: the age-integrated superposition
$F(m) = \sum_t (dN/dm)\,\Delta t$. Open system: the steady state \emph{is}
the instantaneous spectrum, so the last snapshot is the estimator and
$\Delta t$ never enters.

\paragraph{Local slope and the plateau.} The honest way to look at a measured
spectrum is
\begin{equation}
  \Gamma(m) \;=\; \frac{d\log F}{d\log m},
\end{equation}
computed by least squares on a sliding window. A genuine inertial range is a
\emph{plateau} in $\Gamma$; the contaminated ends show up as the places where
$\Gamma$ bends. The reported index is the mean of $\Gamma$ over the longest
contiguous run on which it stays flat within a tolerance, together with its
scatter and its width in decades. A plateau narrower than about one decade is
not a power law, however tight its error bar; the width is reported next to
the index for exactly that reason. If no run qualifies, the answer is
\code{nan}, which is correct rather than a failure.

\paragraph{Guard band.} Independently, strip $0.5$~dex from each end of the
interval between the two characteristic masses --- $[\minj,\msink]$ for an
open system, $[m_i, m_0^{\max}]$ for a closed one --- and fit a single power
law there. This band is chosen \emph{before} looking at the data. The two
estimates must agree; if they do not, the cascade has not developed over the
range that was assumed.

Fitting the whole array does neither. On one closed coagulation run with
$\lambda=0$ and theory $\alpha=-1$: whole array $-1.71$, guard band $-1.09$,
plateau $-1.16 \pm 0.15$ over $1.3$ decades.

\paragraph{Compensated display.} Panels show
\begin{equation}
  m^2 \frac{dN}{dm} \;=\; \text{mass per logarithmic mass interval},
\end{equation}
whose slope is $\alpha+2$. All fitting is done on the raw $dN/dm$; the
compensation is applied at draw time only. The model line carries the plateau
index and an amplitude obtained as the least-squares intercept at
\emph{fixed} slope, so no second fit can quietly move the exponent, and the
theory line uses the same anchoring --- the two curves then differ in slope
alone.

\paragraph{Prediction.} With kernel homogeneity $\lambda$,
\begin{equation}
  \beta =
    \begin{cases}
      \lambda & \text{closed (mass-conserving packet, } n \propto m_0^{-2}),\\[2pt]
      \dfrac{1+\lambda}{2} & \text{open (stationary background, } n \propto m_0^{\alpha}),
    \end{cases}
  \qquad
  \alpha = -(1+\beta),
  \qquad
  b = \frac{1}{1-\beta}.
\end{equation}
Coagulation and fragmentation share $\beta$, $b$ and $\alpha$ exactly; only
the sign of the drift flips, so the law is read in $\tau$ going up and in
$\tau_\ast - \tau$ going down.

\begin{center}
\footnotesize
\begin{tabular}{@{}l l c c c@{}}
\toprule
kernel & $K(m_1,m_2)$ & $\lambda$ &
  $\alpha_{\rm closed}$, $b_{\rm closed}$ &
  $\alpha_{\rm open}$, $b_{\rm open}$ \\
\midrule
constant   & $1$                                & $0$   & $-1$,\ \ $1$    & $-3/2$,\ \ $2$ \\
sum        & $m_1^{1/3}+m_2^{1/3}$              & $1/3$ & $-4/3$,\ \ $3/2$& $-5/3$,\ \ $3$ \\
geometric  & $(m_1^{1/3}+m_2^{1/3})^2$          & $2/3$ & $-5/3$,\ \ $3$  & $-11/6$,\ \ $6$ \\
additive   & $m_1+m_2$                          & $1$   & $-2$,\ \ $\infty$ & $-2$,\ \ $\infty$ \\
product    & $m_1 m_2$                          & $2$   & \multicolumn{2}{c}{above gelation: no cascade} \\
\bottomrule
\end{tabular}
\end{center}

\section{Validation}
\label{sec:num_validation}

Constant kernel, $\lambda = 0$. \emph{plateau} is the automatic index,
\emph{guard} the a-priori band, \emph{dec} the plateau width in decades.

\begin{center}
\small
\begin{tabular}{@{}l r r r r r r@{}}
\toprule
case & $b$ & $b_{\rm th}$ & plateau & $\pm$ & dec & guard \\
\midrule
1 closed coagulation      & $0.990$ & $1$ & $-1.077$ & $0.126$ & $2.70$ & $-1.035$ \\
2 closed fragmentation    & $1.001$ & $1$ & $-0.978$ & $0.141$ & $3.00$ & $-1.056$ \\
3 open coagulation        & $2.001$ & $2$ & $-1.475$ & $0.095$ & $3.60$ & $-1.498$ \\
4 open fragmentation, $f=0$   & --- & $2$ & $-2.047$ & $0.183$ & $0.80$ & $-2.354$ \\
4 open fragmentation, $f=0.3$ & --- & $2$ & $-1.550$ & $0.091$ & $3.40$ & $-1.517$ \\
\bottomrule
\end{tabular}
\end{center}

Case~1 also reproduces the exact Smoluchowski solution for the constant
kernel, $m_0(t) = m_i(1+n_0K_1t/2)$, to $\max|{\rm ratio}-1| = 7.8\times10^{-14}$
--- a direct test of Eq.~\eqref{eq:num_clock}. Mass drift is identically zero
in all four cases.

Two rows are meant to be read as diagnoses rather than measurements. Case~4
at $f=0$ returns $-2.05$ over $0.80$ decades: the index has locked onto $-2$
exactly as Sec.~\ref{sec:trick8} predicts, and the narrow plateau announces
that it is not a measurement. And the growth exponent $b$ is absent for
case~4 because, at this budget, over $99\%$ of live particles still descend
from the initial condition and inherit $t_{\rm born}=0$; the age axis is then
degenerate --- age equals wall-clock time for everyone at once --- and the
isochrones degrade into ``spectrum versus time''. The same fact shows up as
$M_{\rm out}/M_{\rm in} > 1$: the sink is draining the initial reservoir,
whose mass never entered $M_{\rm in}$.

\section{Practical failure modes}
\label{sec:num_pitfalls}

\paragraph{A mass outside the grid is a hard crash.}
The majorant is built from bin \emph{upper corners}, and $b(m)$ clamps
anything above the top edge into the last bin --- whose corner is then
smaller than the particle itself. Hence $K(m,m) > \fmaj_{B-1,B-1}$ and the
first pair raises \code{Majorant violated}. With a constant kernel this can
never happen, since $K\equiv1$ is flat; for any $\lambda>0$ it is immediate.
Check that the grid brackets $\minj$ and $\msink$ before starting, not after.

\paragraph{An unbounded event budget is safe in exactly one case.}
Removing the event cap is only safe when something else can terminate the
run:

\begin{center}
\small
\begin{tabular}{@{}l l p{7.4cm}@{}}
\toprule
case & unbounded & what stops it, or does not \\
\midrule
closed coagulation   & safe   & \code{stop\_max\_mass}; and $N$ decreases monotonically,
                                so \code{live}~$<2$ is a guaranteed terminator \\
closed fragmentation & unsafe & $N$ grows; the stall guard fires only once
                                \emph{everything} sits below the passive floor,
                                i.e.\ at $N_0 m_i/m_{\rm floor}$ particles \\
open coagulation     & hangs  & $N$ sits on a plateau, so no other condition
                                can ever fire \\
open fragmentation   & unsafe & $N$ grows without bound \\
\bottomrule
\end{tabular}
\end{center}

In the open cases the correct brake is \code{stop\_sink\_events}
(Sec.~\ref{sec:num_sink}); keeping a resource ceiling in addition costs
nothing and bounds a mistuned injection rate. The recorded stop reason then
says which one bit.

\paragraph{Changing the kernel invalidates four constants.}
The homogeneity $\lambda$ and every predicted exponent follow automatically,
and the plateau finder, the guard band and the compensation are all
kernel-agnostic. Four things are not:

\begin{enumerate}
  \item \textbf{Injection rate.} Number balance in steady state is
        $q = \langle K\rangle N^2/(2V)$. Writing $q = N^2/2$ is that formula
        at $\langle K\rangle = 1$. Since $N$ is dominated by the injection
        scale when $\alpha<-1$, most pairs are $(\minj,\minj)$, so
        $\langle K \rangle \sim K(\minj,\minj)$, with heavier partners
        pushing it up by roughly a factor two.
  \item \textbf{Collision time.} $t_c = 1/(Kn)$, not $1/n$.
  \item \textbf{Age binning.} One age bin of width $\Delta_\tau$ dex spans
        $b\,\Delta_\tau$ dex of \emph{mass}. Fix the mass resolution and let
        the age step follow, $\Delta_\tau = 0.30/b$; otherwise at $b=6$ a
        step tuned for $b=2$ jumps almost a decade of mass per bin and a
        two-decade fitting window catches two bins.
  \item \textbf{Exact solutions.} The closed-form Smoluchowski reference and
        the ``$m_0(t)$ must be linear'' check are properties of the constant
        kernel alone.
\end{enumerate}

\paragraph{The $-2$ fingerprint has three distinct causes.}
An index of $-2$ is returned by (i) a genuine $\lambda = 1$ kernel, where it
is the correct answer; (ii) the non-local disruption rule of
Sec.~\ref{sec:trick8}, where it is a boundary effect; and (iii) a clock
advanced per event rather than per trial, Sec.~\ref{sec:trick5}, where it is
a numerical artefact. All three drive $\beta \to 1$ and look identical on the
page. Seeing $-2$ therefore says nothing on its own. The discriminating test
is whether changing the kernel changes the index: if it does not, this is not
universality.

\paragraph{A flat shelf is not an inertial range.}
In closed fragmentation the passive floor freezes everything below it: those
particles never fragment again and pile up in a flat shelf that is not part
of the cascade. Left in the array that shelf is the longest flat run there
is --- $4.3$ decades at $\alpha \simeq 0$ --- and the plateau finder duly
reports it. Restrict the search to $m \ge m_{\rm floor}$.

% ---------------------------------------------- delete when including
\end{document}

%      next to the existing \include{Parts/unified-cascade_v4}.
%      main.tex already loads amsmath, amssymb, graphicx and hyperref; you
%      must additionally add to its preamble:
%          \usepackage{algorithm}
%          \usepackage{algpseudocode}
%          \usepackage{booktabs}
% =====================================================================

% ---------------------------------------------- PREAMBLE -- CUT FROM HERE
\documentclass[a4paper,12pt]{article}

\usepackage[T1]{fontenc}
\usepackage[utf8]{inputenc}
\usepackage[english]{babel}
\usepackage{amsmath,amssymb}
\usepackage{geometry}
\geometry{margin=2.5cm}
\usepackage{graphicx}
\usepackage{booktabs}
\usepackage{algorithm}
\usepackage{algpseudocode}
\usepackage[colorlinks=true,linkcolor=blue,citecolor=blue,urlcolor=blue]{hyperref}

\newcommand{\minj}{m_{\mathrm{inj}}}
\newcommand{\msink}{m_{\mathrm{sink}}}
\newcommand{\fmaj}{f^{\mathrm{maj}}}
\newcommand{\code}[1]{\texttt{#1}}

\begin{document}

\title{The numerical scheme: majorant-frequency Monte Carlo\\ for the mass cascade}
\author{}
\date{}
\maketitle
% ---------------------------------------------- PREAMBLE -- CUT TO HERE

\section{Scope and design principle}
\label{sec:num_scope}

The engine integrates the stochastic mass cascade in four configurations,
\[
  \text{process} \in \{\text{coagulation},\ \text{fragmentation}\},
  \qquad
  \text{system} \in \{\text{closed},\ \text{open}\},
\]
through a \emph{single} code path. Only two things differ between the four
cases:

\begin{enumerate}
  \item what an accepted event does to the two selected particles, and
  \item the boundary conditions --- injection plus an absorbing sink, or
        neither.
\end{enumerate}

Everything else --- majorant construction, pair selection, the clock, the
bookkeeping --- is shared. This is a deliberate constraint, not a
convenience: a discrepancy between two of the four cases cannot then be
blamed on two different integrators.

The variant documented here carries no particle weights. One simulated
particle is one physical particle, $w \equiv 1$ for the whole run, so
$\sum_i w_i$ and $\sum_i w_i m_i$ are conserved to machine precision at
every step. The measured mass drift is identically zero, and any nonzero
value is a bug rather than a statistical excursion. Section~\ref{sec:num_cost}
states what that costs.

\subsection{State}

The state is four parallel arrays over live particles, plus a binning of
mass space:

\begin{center}
\small
\begin{tabular}{@{}ll@{}}
\toprule
\code{mass[i]}       & particle mass $m_i$ \\
\code{inj\_time[i]}  & birth time, used for the isochrones (Sec.~\ref{sec:trick7}) \\
\code{bin\_of[i]}    & index of the mass bin holding $i$ \\
\code{pos\_in\_bin[i]} & position of $i$ inside that bin's index list \\
\midrule
\code{edges[0..B]}   & $B+1$ logarithmically spaced mass-bin edges \\
\code{bins[b]}       & list of particle indices in bin $b$ \\
\code{N\_bin[b]}     & occupancy $N_b$ \\
\bottomrule
\end{tabular}
\end{center}

A particle of mass $m$ belongs to bin
$b(m) = \max\bigl(0,\min(B-1,\ \mathrm{searchsorted}(\text{edges},m)-1)\bigr)$.
The clamping in that expression is a trap and is discussed in
Sec.~\ref{sec:num_pitfalls}.

\section{The baseline algorithm}
\label{sec:num_algorithm}

\subsection{Setup}

\begin{algorithm}[h]
\caption{Setup: bins, majorant table, rate vector}
\label{alg:setup}
\begin{algorithmic}[1]
\State bin every initial particle; form the occupancies $N_b$
\State \textbf{majorant table:} for all $b,c$ \quad
       $\fmaj_{bc} \gets K(\text{edges}[b+1],\ \text{edges}[c+1])$
       \Comment{upper corners}
\State $T_b \gets \sum_c \fmaj_{bc} N_c$
\State $\text{row}_b \gets N_b\,(T_b - \fmaj_{bb})$
\State $R \gets \sum_b \text{row}_b$
\end{algorithmic}
\end{algorithm}

The corner rule $\fmaj_{bc}=K(\text{edges}[b+1],\text{edges}[c+1])$ bounds
$K$ over the whole bin pair \emph{provided $K$ is non-decreasing in both
arguments}. That is the single structural assumption on the kernel, and it
is verified at run time rather than assumed: every accepted trial checks
$K(m_i,m_j) \le \fmaj_{bc}$ and raises if it fails.

By construction
\begin{equation}
  R \;=\; \sum_b N_b\bigl(T_b - \fmaj_{bb}\bigr)
    \;=\; \sum_{i \neq j} \fmaj_{ij}
    \;=\; 2 \sum_{i<j} \fmaj_{ij},
  \label{eq:num_R}
\end{equation}
i.e.\ $R$ runs over \emph{ordered} pairs. The factor two in
Eq.~\eqref{eq:num_R} is the origin of the factor two in the clock,
Eq.~\eqref{eq:num_clock}.

\subsection{Main loop}

\begin{algorithm}[h]
\caption{Main loop: one majorant trial}
\label{alg:mainloop}
\begin{algorithmic}[1]
\While{no stopping condition is met}
  \State $\Delta t \gets 2V/(wR)$; \quad $t \gets t + \Delta t$
         \Comment{the clock, Sec.~\ref{sec:trick5}}
  \State \textbf{injection (open systems):}
         $\nu \gets \nu + q\,\Delta t / w$; \quad
         $k \gets \lfloor \nu \rfloor$; \quad $\nu \gets \nu - k$;
         add $k$ particles of mass $\minj$
  \State draw a row bin $b$ with $P(b) \propto \text{row}_b$
  \State draw a column bin $c$ with
         $P(c) \propto N_c \fmaj_{bc}$, using $(N_b-1)\fmaj_{bb}$ for $c=b$
  \State draw $i$ uniformly from \code{bins[b]}, $j \neq i$ uniformly from \code{bins[c]}
  \State \textbf{thin:} accept with probability $K(m_i,m_j)/\fmaj_{bc}$;
         otherwise \textbf{continue} \Comment{null event}
  \State apply the event (Alg.~\ref{alg:event}); apply the absorbing sink
  \State record a snapshot if the cadence says so (Sec.~\ref{sec:trick10})
\EndWhile
\end{algorithmic}
\end{algorithm}

A rejected trial is not an error: it is the thinning that makes the
dominating process exact. The acceptance ratio, reported as \code{acc} in
the run log, is $O(1)$ by construction of the bin-pair table --- measured
$0.95$ for the geometric kernel on bins of $0.1$~dex, where the corner
overshoot is $10^{0.1\lambda}$ per argument.

\subsection{The event}

\begin{algorithm}[h]
\caption{The accepted event --- the only place the four cases differ}
\label{alg:event}
\begin{algorithmic}[1]
\If{coagulation}
  \State $m_{\mathrm{new}} \gets m_i + m_j$
  \State inherit the age of $i$ from $(i,j)$ by \code{age\_rule} (Sec.~\ref{sec:trick7})
  \State $m_i \gets m_{\mathrm{new}}$; rebin $i$; \textbf{delete} $j$
  \State \textbf{if} open \textbf{and} $m_{\mathrm{new}} \ge \msink$ \textbf{then}
         delete $i$; $M_{\mathrm{out}} \mathrel{+}= w\,m_{\mathrm{new}}$;
         $n_{\mathrm{out}} \mathrel{+}= 1$
\Else \Comment{fragmentation}
  \State $p \gets \arg\max(m_i,m_j)$; \quad $s \gets \arg\min(m_i,m_j)$
  \State \textbf{if} $m_s < f\,m_p$ \textbf{then continue}
         \Comment{disruption threshold, Sec.~\ref{sec:trick8}}
  \State $\xi \sim U(0,1)$; \quad $m_a \gets \xi m_p$, \ $m_b \gets (1-\xi)m_p$
  \State both fragments inherit $t_{\mathrm{born}}$ of the single parent $p$
  \State $m_p \gets m_a$; rebin $p$; \textbf{append} a particle of mass $m_b$
  \State \textbf{if} open: delete every fragment with $m \le \msink$,
         highest index first
\EndIf
\end{algorithmic}
\end{algorithm}

Note the sign asymmetry, which governs everything in
Sec.~\ref{sec:num_pitfalls}: coagulation is $2 \to 1$ and therefore a
\emph{sink} of particle number, fragmentation is $1 \to 2$ and therefore a
\emph{source}.

\section{The ten tricks}
\label{sec:num_tricks}

\subsection{Majorant collision frequency}
\label{sec:trick1}

Computing the exact pair rate for every pair is $O(N^2)$ per step. Instead a
majorant $\fmaj \ge K$ drives a dominating Poisson process and events are
thinned with acceptance $K/\fmaj$. The thinned process is exact: the
majorant affects the cost, never the statistics.

\emph{Sources.} Ivanov \& Rogasinsky, Sov.\ J.\ Numer.\ Anal.\ Math.\
Modelling \textbf{3}, 453 (1988) [MCF scheme in DSMC]; Garcia, van den
Broeck, Aertsens \& Serneels, Physica~A \textbf{143}, 535 (1987) [majorant
for coagulation]; Gillespie, J.~Atmos.\ Sci.\ \textbf{32}, 1977 (1975)
[exact stochastic coalescence].

\subsection{Bin-pair majorant table}
\label{sec:trick2}

The majorant is a matrix $\fmaj_{bc}$ over pairs of mass bins, not a single
global number. What that buys is \emph{cost} and \emph{checkability}. It is
worth being precise about which, because a majorant cannot buy correctness:
it cancels out of the physical rate exactly.

\paragraph{The majorant cancels.}
Ticks of the dominating process arrive at rate $wR/(2V)$; a tick selects the
ordered pair $(i,j)$ with probability $\fmaj_{ij}/R$; it is accepted with
probability $K_{ij}/\fmaj_{ij}$. The rate of accepted events on that pair is
therefore
\begin{equation}
  \frac{wR}{2V}\cdot\frac{\fmaj_{ij}}{R}\cdot\frac{K_{ij}}{\fmaj_{ij}}
  \;=\; \frac{w\,K_{ij}}{2V},
  \label{eq:num_cancel}
\end{equation}
and summing over unordered pairs gives $w\sum_{i<j}K_{ij}/V$ --- the
Smoluchowski rate, independent of $\fmaj$ altogether. The majorant cancels
twice, once in the selection weight and once in the acceptance. \emph{Any}
valid majorant therefore returns the same physics; a bad one only wastes
trials. Equation~\eqref{eq:num_cancel} holds because the clock is advanced on
every \emph{trial} rather than on every event --- see
Sec.~\ref{sec:trick5}, which spells out the one edit that breaks it.

\paragraph{What the table actually buys.}
A single global $f_{\max} = K(m_{\max},m_{\max})$ is a legitimate majorant and
gives identical statistics. It is nevertheless the wrong choice, for three
reasons:

\begin{itemize}
  \item \textbf{Acceptance.} A typical pair is accepted with probability
        $\langle K\rangle/f_{\max} \simeq (m_0/m_{\max})^{\lambda}$. Across
        $D$ decades that is $10^{-\lambda D}$: for the geometric kernel over
        four decades, $2\times10^{-3}$, i.e.\ some five hundred trials per
        event against $1.05$ with the table. The table instead keeps the
        acceptance $O(1)$ in \emph{every} bin pair, because the bound is
        evaluated at the corners of the pair actually drawn; the overshoot is
        then only $10^{0.1\lambda}$ per argument at $0.1$~dex bins.
  \item \textbf{The stall guard misfires.} \code{max\_stall\_tries} counts
        trials without an accepted event. At an acceptance of
        $2\times10^{-3}$ its default of $2\times10^{6}$ is only
        $\sim\!4000$ events' worth of honest work, so a healthy run can be
        killed with the diagnosis ``acceptance collapsed''.
  \item \textbf{A lagging maximum is a genuine bias.} A \emph{running} global
        maximum that is not refreshed before some particle grows past it
        gives $\fmaj < K$, and the thinning is then silently wrong --- this,
        unlike the two points above, is a correctness failure. The corner
        table is static, and every trial checks $K \le \fmaj$ explicitly, so
        the same situation raises instead of biasing.
\end{itemize}

\emph{Source.} Eibeck \& Wagner, SIAM J.~Sci.\ Comput.\ \textbf{22}, 802
(2000),
\url{https://doi.org/10.1137/S1064827599353488}.

\subsection{$O(1)$ bin membership}
\label{sec:trick3}

\code{bins[b]} is a list of particle indices and \code{pos\_in\_bin[i]} is
the position of $i$ inside it. Removal swaps the last element into the hole,
so insertion and deletion are $O(1)$ with no reallocation. The same
swap-and-pop idiom is used on the particle arrays themselves when a particle
is deleted, which is why every index held across a deletion must be
revalidated.

\subsection{Incremental rate maintenance}
\label{sec:trick4}

When one bin count changes by $\delta$,
\[
  T \mathrel{+}= \delta\, \fmaj_{\cdot b},
  \qquad
  \text{row} \gets N \odot (T - \mathrm{diag}\,\fmaj),
  \qquad
  R \gets \textstyle\sum_b \text{row}_b ,
\]
an $O(B)$ vector operation instead of recomputing the full matrix--vector
product. Selection uses a vectorised \code{cumsum} plus \code{searchsorted}
rather than a Python loop; that loop is otherwise the dominant cost. For
very large $B$ a Fenwick tree gives $O(\log B)$ (Goodson \& Kraft,
J.~Comput.\ Phys.\ \textbf{183}, 210 (2002)).

\subsection{The clock}
\label{sec:trick5}

The step is
\begin{equation}
  \Delta t \;=\; \frac{2V}{w\,R},
  \label{eq:num_clock}
\end{equation}
and three separate things about it are easy to get wrong.

\paragraph{Ordered versus unordered pairs.}
$R$ counts every unordered pair twice, Eq.~\eqref{eq:num_R}, while the
physical majorant rate runs over unordered pairs. Hence the numerator is
$2V$, not $V$. Pair \emph{selection} is unaffected --- both orderings lead to
the same unordered pair and the factors cancel --- so the error is invisible
in the dynamics and rescales $t$ by exactly two, breaking every comparison
with an analytic solution.

\paragraph{Weights.}
Here $w$ is a constant, but it is still written into the step so that the
formula matches the weighted variant exactly and a constant $w \neq 1$ simply
rescales the density. Nothing in this variant ever changes $w$, so the clock
cannot be corrupted by weight bookkeeping. That is the entire reason the
baseline exists.

\paragraph{Quadratic scaling.}
$R \propto N^2$ because the process is binary, so
\begin{equation}
  t \;\simeq\; \frac{2E}{N^2}
  \label{eq:num_t_of_E}
\end{equation}
after $E$ trials. A step linear in $N$ describes a unary process and gives
the wrong $N(t)$. Equation~\eqref{eq:num_t_of_E} has a practical corollary
that is easy to meet by accident: \emph{raising $N$ at fixed event budget
makes the run shorter in physical time, not longer}. Any gate expressed as an
absolute time and tuned at one $N$ is never crossed at $10N$.

\paragraph{Time advances on every trial, not on every event.}
\code{t\_phys} is incremented immediately after $\Delta t$ is formed, some
fifty lines before the acceptance test, so null (thinned) trials move the
clock too. That is precisely what makes the cancellation
\eqref{eq:num_cancel} work, and moving the increment inside the accepted
branch is the one edit that turns the majorant from a cost into a
\emph{bias}. Physical time would then advance by $2V/(wR)$ per accepted
event, so the reported event rate would be $R$ instead of
$R\,\langle K/\fmaj\rangle = \sum K$, and with a global $f_{\max}$,
\[
  \Delta t_{\rm eff} \;\propto\; \frac{1}{N^2 f_{\max}}
                     \;\propto\; \frac{1}{N^2 m_{\max}^{\lambda}} .
\]
The kernel would then enter through the \emph{extreme} of the distribution
instead of through the typical mass. In an open system $m_{\max}$ is pinned
at the sink and does not grow, so the $m_0$ dependence leaves the clock
entirely, the drift degenerates to $dm_0/dt \propto m_0$ (exponential,
$\beta = 1$), and the constant-flux argument then returns $\alpha = -2$ for
every kernel. This is a third, purely numerical route to the same $-2$
fingerprint that Sec.~\ref{sec:trick8} produces physically.

\subsection{Deterministic residual injection}
\label{sec:trick6}

Injection uses a fractional accumulator,
\[
  \nu \mathrel{+}= q\,\Delta t/w,
  \qquad k = \lfloor \nu \rfloor,
  \qquad \nu \mathrel{-}= k,
\]
rather than a Poisson draw. This removes shot noise at the injection scale.
It is a modelling choice, and it is recorded in the output metadata rather
than left implicit.

\subsection{Isochrone bookkeeping}
\label{sec:trick7}

Every particle carries its birth time; its age is $\tau = t - t_{\rm born}$.
Snapshots accumulate a two-dimensional histogram in $(\tau, m)$, whose rows
are the isochrones and whose sum over rows is the time-averaged spectrum.

For \textbf{fragmentation} the age of a fragment is unambiguous: fragments
descend from exactly one parent. For \textbf{coagulation} a merged particle
has two ancestors and an inheritance rule must be chosen. It is exposed as a
parameter, \code{age\_rule}, precisely because the isochrone analysis depends
on it and the choice must be reported and tested:
\[
  t_{\rm born} \;=\;
  \begin{cases}
    \min(t_1,t_2)                    & \text{\code{'min'}: the oldest ancestor wins,}\\[2pt]
    t_1 \ \text{if}\ m_1 \ge m_2     & \text{\code{'heavier'},}\\[2pt]
    \dfrac{m_1 t_1 + m_2 t_2}{m_1+m_2} & \text{\code{'mass\_weighted'}.}
  \end{cases}
\]

\subsection{Locality of the fragmentation rule --- a physics trap}
\label{sec:trick8}

If the rule is ``the heavier of the pair breaks, whatever hit it'', the
disruption rate of a body of mass $m_0$ against a power-law background
$F(m) \propto m^{\alpha}$ is
\begin{equation}
  \nu(m_0) \;\sim\; \int_{\msink}^{m_0} F(m')\,K(m_0,m')\,dm'
           \;\sim\; \int_{\msink/m_0} x^{\alpha}\,dx ,
  \qquad x = m'/m_0 ,
\end{equation}
which \emph{diverges at the lower limit whenever $\alpha < -1$}. The drift is
then set by the smallest particles in the box --- by the sink scale, not by
$m_0$. The mean-field closure is void, the drift degenerates to
$dm_0/d\tau \sim -C m_0$ (i.e.\ $\beta = 1$, exponential decay), and the
measured index locks onto $-2$ regardless of the kernel: a spurious
universality that is really a boundary effect.

The cure is the one Dohnanyi (1969) built in --- a catastrophic-disruption
threshold. A grain of dust does not shatter a boulder. Requiring
\begin{equation}
  m_{\rm small} \;\ge\; f\,m_{\rm large},
  \qquad f = O(0.1\text{--}1),
\end{equation}
restricts the integral to $x \in [f,1]$, which is scale free, restores
locality, and recovers $\alpha = -(3+\lambda)/2$.

A \emph{closed} system is immune: its self-similar packet is narrow, so all
partners are already of order $m_0$ and the rule is local by construction.
This is why the closed runs agree with the theory at $f=0$ and the open ones
do not.

\subsection{No resampling --- and what it costs}
\label{sec:num_cost}

What is bought:
\begin{itemize}
  \item $\sum w$ and $\sum wm$ conserved to machine precision; the reported
        \code{mass\_drift} is identically zero and any nonzero value is a
        genuine bug;
  \item no clock subtlety at all, $\Delta t = 2V/R$ with $w=1$;
  \item nothing can bias the result, because nothing is resampled.
\end{itemize}

What is paid, and it is the binding constraint on every run:
\begin{itemize}
  \item \textbf{Coagulation.} In a closed box $N = N_i m_i/m_0$, so keeping
        $\sim100$ particles alive at the top requires $m_0 \lesssim N_i/100$:
        $N_i = 10^6$ buys four decades, $N_i = 10^8$ buys six (and
        $\sim3$~GB of arrays). Pushed past that, a run ends with a handful of
        particles; at $N_i = 2\times10^5$ one test finished with a single
        particle alive and the usable plateau shrank from $5.6$ decades to
        $2.7$.
  \item \textbf{Fragmentation.} Worse, in the other direction: $N$ grows as
        $m_i/m_0$, so a completed cascade holds $N_0\,\minj/\msink$
        particles. Three decades multiply the population a thousandfold, and
        \emph{memory, not physics, ends the run}.
\end{itemize}

Use this variant when the priority is an unimpeachable baseline. Use the
weighted variant, which adds multiplication and down-sampling, when the
priority is range: ten decades there cost $4.5\times10^5$ events and returned
$\alpha = -1.005$ against a theoretical $-1.000$. The two agree wherever both
can run, which is the point of keeping both.

\subsection{Snapshot placement}
\label{sec:trick10}

The closed-system observable is a Riemann sum
$F(m) = \sum_k (dN/dm)(m,t_k)\,\Delta t_k$, and the sampling of $t$ has two
separate jobs that must not be confused. \textbf{Placement} decides what is
\emph{resolved}: for a given $m$ the integrand peaks at $m_0 \sim m$, so
every decade of $m$ needs snapshots at the matching decade of $m_0$.
\textbf{Weight} decides what is \emph{computed}: it must be the actual
$\Delta t$ between consecutive snapshots, whatever the placement.

\begin{center}
\small
\begin{tabular}{@{}l p{7.2cm} l@{}}
\toprule
placement & what happens & measured $\alpha$ (theory $-1$) \\
\midrule
\code{events} & events per unit $m_0$ fall as $m_0^{-2}$, so $\Delta t$
  explodes; $94\%$ of the total weight lands in one snapshot
  & $-1.4 \ldots -2.6$ \\
\code{t\_phys}, narrow & fine over $\lesssim 2$ decades & $-0.97 \ldots -1.02$ \\
\code{t\_phys}, wide & $t \propto m_0^{1-\lambda}$, so uniform $t$ is nearly
  uniform $m_0$; over ten decades the lower five get no sample at all
  & $-0.04 \pm 0.16$ \\
\code{log\_m0} & $k$ samples per decade of $m_0$, everywhere
  & $\mathbf{-1.005}$, $R^2 = 0.9995$ \\
\bottomrule
\end{tabular}
\end{center}

The failure of \code{events} deserves emphasis, because it is silent: with
that placement the plateau finder of Sec.~\ref{sec:num_measure} still reports
a clean one-decade plateau at $-1.375 \pm 0.14$. An internally consistent,
entirely wrong measurement. \emph{Fitting the plateau cannot rescue a broken
estimator}, because it tests whether the curve you computed has a scaling
region, not whether that curve is the right quantity. The independent guard
is the printed weight of the last snapshot: it must be $\ll 1\%$.

Resampling (Sec.~\ref{sec:num_cost}) and placement cure different diseases.
Resampling buys the \emph{range}; placement buys a \emph{correct
measurement} over that range.

\paragraph{Rule.} Closed system $\Rightarrow$ \code{log\_m0}, stride $=$
snapshots per decade ($20$--$30$). Open system $\Rightarrow$ the steady state
is the \emph{last} snapshot, $\Delta t$ never enters the estimator, and any
cadence works.

\section{The sink as a criterion}
\label{sec:num_sink}

A stopping or gating criterion phrased in events is a statement about
\emph{resources}; one phrased in sink absorptions is a statement about the
\emph{physics}. Only the second survives a change of $N$, of mass range, or
of kernel. The engine therefore counts $n_{\rm out}$, the number of physical
particles that have left through the absorbing boundary, and exposes it
three ways.

\begin{center}
\small
\begin{tabular}{@{}l p{8.2cm}@{}}
\toprule
\code{iso\_start\_sink} (default $10$) &
  begin accumulating isochrones only after $n_{\rm out}$ absorptions \\
\code{stop\_sink\_events} (default \code{None}) &
  stop the run at $n_{\rm out}$ absorptions \\
\code{out["n\_out"]} &
  the counter as a per-snapshot time series \\
\bottomrule
\end{tabular}
\end{center}

The gate says: bin no age until the cascade has demonstrably pushed that
many particles across the \emph{whole} inertial range. Its cost follows from
the mass budget --- the box holds $\simeq N\sqrt{\msink}$ in steady state and
$n_{\rm out}$ absorptions carry off $n_{\rm out}\msink$ --- so
\begin{equation}
  E \;\simeq\; N\sqrt{\msink} \;+\; n_{\rm out}\,\msink .
  \label{eq:num_budget}
\end{equation}
The second term is independent of $N$: the gate does not get more expensive
as the run grows, only filling the box does. Equation~\eqref{eq:num_budget}
is derived at $\lambda = 0$ and overestimates the cost for $\lambda > 0$,
where a steeper kernel drives mass to the sink faster per event.

A closed system has no sink, so the armed default is silently disarmed
there; only an explicitly passed value raises. The distinction between ``the
caller typed this'' and ``this came from the signature'' is carried by an
\code{int} subclass, so the value behaves as a plain integer everywhere else.

\section{The measurement layer}
\label{sec:num_measure}

A run produces arrays; turning them into an exponent is a separate stage
with its own failure modes.

\paragraph{Estimator.} Closed system: the age-integrated superposition
$F(m) = \sum_t (dN/dm)\,\Delta t$. Open system: the steady state \emph{is}
the instantaneous spectrum, so the last snapshot is the estimator and
$\Delta t$ never enters.

\paragraph{Local slope and the plateau.} The honest way to look at a measured
spectrum is
\begin{equation}
  \Gamma(m) \;=\; \frac{d\log F}{d\log m},
\end{equation}
computed by least squares on a sliding window. A genuine inertial range is a
\emph{plateau} in $\Gamma$; the contaminated ends show up as the places where
$\Gamma$ bends. The reported index is the mean of $\Gamma$ over the longest
contiguous run on which it stays flat within a tolerance, together with its
scatter and its width in decades. A plateau narrower than about one decade is
not a power law, however tight its error bar; the width is reported next to
the index for exactly that reason. If no run qualifies, the answer is
\code{nan}, which is correct rather than a failure.

\paragraph{Guard band.} Independently, strip $0.5$~dex from each end of the
interval between the two characteristic masses --- $[\minj,\msink]$ for an
open system, $[m_i, m_0^{\max}]$ for a closed one --- and fit a single power
law there. This band is chosen \emph{before} looking at the data. The two
estimates must agree; if they do not, the cascade has not developed over the
range that was assumed.

Fitting the whole array does neither. On one closed coagulation run with
$\lambda=0$ and theory $\alpha=-1$: whole array $-1.71$, guard band $-1.09$,
plateau $-1.16 \pm 0.15$ over $1.3$ decades.

\paragraph{Compensated display.} Panels show
\begin{equation}
  m^2 \frac{dN}{dm} \;=\; \text{mass per logarithmic mass interval},
\end{equation}
whose slope is $\alpha+2$. All fitting is done on the raw $dN/dm$; the
compensation is applied at draw time only. The model line carries the plateau
index and an amplitude obtained as the least-squares intercept at
\emph{fixed} slope, so no second fit can quietly move the exponent, and the
theory line uses the same anchoring --- the two curves then differ in slope
alone.

\paragraph{Prediction.} With kernel homogeneity $\lambda$,
\begin{equation}
  \beta =
    \begin{cases}
      \lambda & \text{closed (mass-conserving packet, } n \propto m_0^{-2}),\\[2pt]
      \dfrac{1+\lambda}{2} & \text{open (stationary background, } n \propto m_0^{\alpha}),
    \end{cases}
  \qquad
  \alpha = -(1+\beta),
  \qquad
  b = \frac{1}{1-\beta}.
\end{equation}
Coagulation and fragmentation share $\beta$, $b$ and $\alpha$ exactly; only
the sign of the drift flips, so the law is read in $\tau$ going up and in
$\tau_\ast - \tau$ going down.

\begin{center}
\footnotesize
\begin{tabular}{@{}l l c c c@{}}
\toprule
kernel & $K(m_1,m_2)$ & $\lambda$ &
  $\alpha_{\rm closed}$, $b_{\rm closed}$ &
  $\alpha_{\rm open}$, $b_{\rm open}$ \\
\midrule
constant   & $1$                                & $0$   & $-1$,\ \ $1$    & $-3/2$,\ \ $2$ \\
sum        & $m_1^{1/3}+m_2^{1/3}$              & $1/3$ & $-4/3$,\ \ $3/2$& $-5/3$,\ \ $3$ \\
geometric  & $(m_1^{1/3}+m_2^{1/3})^2$          & $2/3$ & $-5/3$,\ \ $3$  & $-11/6$,\ \ $6$ \\
additive   & $m_1+m_2$                          & $1$   & $-2$,\ \ $\infty$ & $-2$,\ \ $\infty$ \\
product    & $m_1 m_2$                          & $2$   & \multicolumn{2}{c}{above gelation: no cascade} \\
\bottomrule
\end{tabular}
\end{center}

\section{Validation}
\label{sec:num_validation}

Constant kernel, $\lambda = 0$. \emph{plateau} is the automatic index,
\emph{guard} the a-priori band, \emph{dec} the plateau width in decades.

\begin{center}
\small
\begin{tabular}{@{}l r r r r r r@{}}
\toprule
case & $b$ & $b_{\rm th}$ & plateau & $\pm$ & dec & guard \\
\midrule
1 closed coagulation      & $0.990$ & $1$ & $-1.077$ & $0.126$ & $2.70$ & $-1.035$ \\
2 closed fragmentation    & $1.001$ & $1$ & $-0.978$ & $0.141$ & $3.00$ & $-1.056$ \\
3 open coagulation        & $2.001$ & $2$ & $-1.475$ & $0.095$ & $3.60$ & $-1.498$ \\
4 open fragmentation, $f=0$   & --- & $2$ & $-2.047$ & $0.183$ & $0.80$ & $-2.354$ \\
4 open fragmentation, $f=0.3$ & --- & $2$ & $-1.550$ & $0.091$ & $3.40$ & $-1.517$ \\
\bottomrule
\end{tabular}
\end{center}

Case~1 also reproduces the exact Smoluchowski solution for the constant
kernel, $m_0(t) = m_i(1+n_0K_1t/2)$, to $\max|{\rm ratio}-1| = 7.8\times10^{-14}$
--- a direct test of Eq.~\eqref{eq:num_clock}. Mass drift is identically zero
in all four cases.

Two rows are meant to be read as diagnoses rather than measurements. Case~4
at $f=0$ returns $-2.05$ over $0.80$ decades: the index has locked onto $-2$
exactly as Sec.~\ref{sec:trick8} predicts, and the narrow plateau announces
that it is not a measurement. And the growth exponent $b$ is absent for
case~4 because, at this budget, over $99\%$ of live particles still descend
from the initial condition and inherit $t_{\rm born}=0$; the age axis is then
degenerate --- age equals wall-clock time for everyone at once --- and the
isochrones degrade into ``spectrum versus time''. The same fact shows up as
$M_{\rm out}/M_{\rm in} > 1$: the sink is draining the initial reservoir,
whose mass never entered $M_{\rm in}$.

\section{Practical failure modes}
\label{sec:num_pitfalls}

\paragraph{A mass outside the grid is a hard crash.}
The majorant is built from bin \emph{upper corners}, and $b(m)$ clamps
anything above the top edge into the last bin --- whose corner is then
smaller than the particle itself. Hence $K(m,m) > \fmaj_{B-1,B-1}$ and the
first pair raises \code{Majorant violated}. With a constant kernel this can
never happen, since $K\equiv1$ is flat; for any $\lambda>0$ it is immediate.
Check that the grid brackets $\minj$ and $\msink$ before starting, not after.

\paragraph{An unbounded event budget is safe in exactly one case.}
Removing the event cap is only safe when something else can terminate the
run:

\begin{center}
\small
\begin{tabular}{@{}l l p{7.4cm}@{}}
\toprule
case & unbounded & what stops it, or does not \\
\midrule
closed coagulation   & safe   & \code{stop\_max\_mass}; and $N$ decreases monotonically,
                                so \code{live}~$<2$ is a guaranteed terminator \\
closed fragmentation & unsafe & $N$ grows; the stall guard fires only once
                                \emph{everything} sits below the passive floor,
                                i.e.\ at $N_0 m_i/m_{\rm floor}$ particles \\
open coagulation     & hangs  & $N$ sits on a plateau, so no other condition
                                can ever fire \\
open fragmentation   & unsafe & $N$ grows without bound \\
\bottomrule
\end{tabular}
\end{center}

In the open cases the correct brake is \code{stop\_sink\_events}
(Sec.~\ref{sec:num_sink}); keeping a resource ceiling in addition costs
nothing and bounds a mistuned injection rate. The recorded stop reason then
says which one bit.

\paragraph{Changing the kernel invalidates four constants.}
The homogeneity $\lambda$ and every predicted exponent follow automatically,
and the plateau finder, the guard band and the compensation are all
kernel-agnostic. Four things are not:

\begin{enumerate}
  \item \textbf{Injection rate.} Number balance in steady state is
        $q = \langle K\rangle N^2/(2V)$. Writing $q = N^2/2$ is that formula
        at $\langle K\rangle = 1$. Since $N$ is dominated by the injection
        scale when $\alpha<-1$, most pairs are $(\minj,\minj)$, so
        $\langle K \rangle \sim K(\minj,\minj)$, with heavier partners
        pushing it up by roughly a factor two.
  \item \textbf{Collision time.} $t_c = 1/(Kn)$, not $1/n$.
  \item \textbf{Age binning.} One age bin of width $\Delta_\tau$ dex spans
        $b\,\Delta_\tau$ dex of \emph{mass}. Fix the mass resolution and let
        the age step follow, $\Delta_\tau = 0.30/b$; otherwise at $b=6$ a
        step tuned for $b=2$ jumps almost a decade of mass per bin and a
        two-decade fitting window catches two bins.
  \item \textbf{Exact solutions.} The closed-form Smoluchowski reference and
        the ``$m_0(t)$ must be linear'' check are properties of the constant
        kernel alone.
\end{enumerate}

\paragraph{The $-2$ fingerprint has three distinct causes.}
An index of $-2$ is returned by (i) a genuine $\lambda = 1$ kernel, where it
is the correct answer; (ii) the non-local disruption rule of
Sec.~\ref{sec:trick8}, where it is a boundary effect; and (iii) a clock
advanced per event rather than per trial, Sec.~\ref{sec:trick5}, where it is
a numerical artefact. All three drive $\beta \to 1$ and look identical on the
page. Seeing $-2$ therefore says nothing on its own. The discriminating test
is whether changing the kernel changes the index: if it does not, this is not
universality.

\paragraph{A flat shelf is not an inertial range.}
In closed fragmentation the passive floor freezes everything below it: those
particles never fragment again and pile up in a flat shelf that is not part
of the cascade. Left in the array that shelf is the longest flat run there
is --- $4.3$ decades at $\alpha \simeq 0$ --- and the plateau finder duly
reports it. Restrict the search to $m \ge m_{\rm floor}$.

% ---------------------------------------------- delete when including
\end{document}
